\chapter{از بازنمایی انتزاعی تا سیگنال فیزیکی}
\label{ch:representation-to-physical-fa}

در پایان فصل قبل به یک پرسش رسیدیم: اگر هر لایه از بازنمایی به لایه‌ای پایین‌تر اشاره کند، این زنجیره سرانجام کجا متوقف می‌شود؟ یک کلمه می‌تواند نماینده‌ی یک مفهوم باشد؛ آن مفهوم می‌تواند چیزی را در ادراک ما بازنمایی کند؛ و ادراک نیز با فعالیت دستگاه عصبی ارتباط داشته باشد. اما اگر بخواهیم همین بازنمایی را از ذهن خود بیرون بیاوریم، روی کاغذ نگه داریم، برای شخص دیگری بفرستیم یا در یک سامانه‌ی محاسباتی از آن استفاده کنیم، دیگر نمی‌توانیم فقط درباره‌ی رابطه‌ی میان دو مفهوم حرف بزنیم. باید چیزی وجود داشته باشد که بتوان آن را دید، شنید، اندازه گرفت یا به هر روش دیگری تشخیص داد.

شاید در نگاه اول پاسخ ساده به نظر برسد: کافی است آن چیز را روی کاغذ بنویسیم یا با صدای بلند بگوییم. اما همین پاسخ ساده چند سؤال دیگر به همراه دارد. چرا جوهر روی کاغذ می‌تواند یک کلمه را نمایندگی کند؟ چرا تغییر فشار هوا می‌تواند جمله‌ای را به گوش شخص دیگری برساند؟ چرا دو وضعیت متفاوت یک وسیله می‌توانند برای ما دو پیام متفاوت داشته باشند؟ و مهم‌تر از همه، چه چیزی باعث می‌شود یک تفاوت فیزیکی را «نشانه» یا «سیگنال» بدانیم، نه صرفاً تفاوتی بی‌اهمیت در جهان؟

برای پاسخ‌دادن به این پرسش‌ها، باید از دو سوی مسئله حرکت کنیم. از یک طرف، باید بفهمیم اطلاعات را چگونه می‌توان توصیف و اندازه‌گیری کرد؛ از طرف دیگر، باید ببینیم آن چیزی که اندازه می‌گیریم در جهان واقعی چگونه حضور پیدا می‌کند. در مسیر نخست با نظریه‌ی اطلاعات، احتمال، آنتروپی، نویز و ظرفیت کانال روبه‌رو می‌شویم. در مسیر دوم، از نشانه و بازنمایی به حالت‌های فیزیکی، سیگنال و حامل می‌رسیم. این دو مسیر در نهایت به یک نقطه‌ی مشترک ختم می‌شوند: اطلاعات را می‌توان به‌صورت انتزاعی توصیف کرد، اما هر بار که آن را ذخیره یا منتقل می‌کنیم، باید از تفاوتی واقعی در جهان فیزیکی استفاده کنیم.

\section{نظریه‌ی اطلاعات: ارتباط را چگونه اندازه بگیریم؟}
\label{sec:information-theory}

فرض کنید می‌خواهید برای دوستی در شهری دیگر پیامی بفرستید. می‌توانید یک نامه بنویسید، با او تماس بگیرید، پیامک بفرستید یا حتی از طریق یک چراغ، علامتی از پیش تعیین‌شده را به او نشان دهید. شکل ارتباط تغییر می‌کند، اما در تمام این حالت‌ها یک مسئله‌ی مشترک وجود دارد: شما چیزی را در اختیار دارید و می‌خواهید شخص دیگری بتواند از روی چیزی که دریافت می‌کند، آن را بازسازی یا تشخیص دهد.

در ساده‌ترین حالت، شاید فقط بخواهید یکی از دو وضعیت را منتقل کنید: «بله» یا «خیر». اما حالا تصور کنید قرار است یکی از صدها واژه، یکی از هزاران تصویر یا یک فایل چند گیگابایتی را منتقل کنید. آیا روش اندازه‌گیری ما برای این پیام‌ها باید کاملاً متفاوت باشد؟ آیا می‌توانیم راهی پیدا کنیم که بدون توجه به اینکه پیام یک جمله است، یک تصویر است یا مجموعه‌ای از اعداد، درباره‌ی مقدار اطلاعات و محدودیت انتقال آن صحبت کنیم؟

این همان مسئله‌ای است که نظریه‌ی اطلاعات به آن می‌پردازد. نظریه‌ی اطلاعات در شکل کلاسیک خود، بیش از آنکه درباره‌ی معنای انسانی پیام‌ها باشد، درباره‌ی امکان بازنمایی، کدگذاری و انتقال آن‌هاست. این تمایز را از ابتدا باید جدی بگیریم؛ چون اگر آن را فراموش کنیم، ممکن است هر عددی که در نظریه‌ی اطلاعات می‌بینیم با مقدار دانش، معنا یا اهمیت یک پیام اشتباه گرفته شود.

\subsection{مدل ارتباط: منبع، رمزگذار، کانال، رمزگشا و مقصد}
\label{sub:communication-model}

کلود شانون در مقاله‌ی مشهور خود، \lr{A Mathematical Theory of Communication}، مسئله‌ی ارتباط را با یک مدل نسبتاً ساده صورت‌بندی کرد. در این مدل، یک منبع اطلاعاتی پیامی را انتخاب می‌کند؛ فرستنده یا رمزگذار آن پیام را به شکلی مناسب برای انتقال تبدیل می‌کند؛ کانال آن را از یک نقطه به نقطه‌ی دیگر می‌رساند؛ گیرنده یا رمزگشا سیگنال دریافت‌شده را تفسیر می‌کند؛ و در پایان، مقصد پیام را دریافت می‌کند. در کنار این مسیر، منبعی از نویز نیز می‌تواند سیگنال را تغییر دهد. \autocite{shannon1948}

برای روشن‌تر شدن ماجرا، فرض کنید دو نفر با بی‌سیم با یکدیگر صحبت می‌کنند. شخص اول جمله‌ای را در ذهن دارد. این جمله منبع پیام نیست، بلکه خودِ پیام در یک صورت زبانی است که قرار است منتقل شود. میکروفون صدا را دریافت می‌کند و فرستنده آن را به سیگنالی مناسب برای ارسال تبدیل می‌کند. موج الکترومغناطیسی از فضای میان دو دستگاه عبور می‌کند. دستگاه دوم آن موج را دریافت می‌کند، سیگنال را از آن استخراج می‌کند و بلندگو دوباره تغییرات الکتریکی را به تغییرات فشار هوا تبدیل می‌کند. در پایان، گوش و مغز شخص دوم صدا را دریافت و تفسیر می‌کنند.

در این مثال چند چیز متفاوت داریم که در گفت‌وگوی روزمره ممکن است همه را «پیام» بنامیم: جمله‌ای که گوینده قصد بیانش را دارد، صدای تولیدشده، سیگنال الکتریکی میکروفون، موج رادیویی، سیگنال بازسازی‌شده در دستگاه مقصد و صدایی که از بلندگو خارج می‌شود. این‌ها یک چیز واحد نیستند؛ بلکه نمایش‌های متفاوتی از یک پیام در مسیر ارتباط‌اند.

\begin{figure}[htbp]
\centering
\begin{tikzpicture}[
  box/.style={draw=genixblue!75,fill=genixblue!5,rounded corners=2pt,minimum width=1.65cm,minimum height=0.9cm,align=center,font=\small},
  arrow/.style={-{Latex[length=2mm]},thick,genixgray},
  node distance=0.48cm
]
  \node[box] (src) {منبع\\پیام};
  \node[box,right=of src] (enc) {رمزگذار\\فرستنده};
  \node[box,right=of enc] (ch) {کانال};
  \node[box,right=of ch] (dec) {رمزگشا\\گیرنده};
  \node[box,right=of dec] (dst) {مقصد};
  \draw[arrow] (src)--node[above,font=\scriptsize]{پیام}(enc);
  \draw[arrow] (enc)--node[above,font=\scriptsize]{سیگنال}(ch);
  \draw[arrow] (ch)--node[above,font=\scriptsize]{سیگنال دریافتی}(dec);
  \draw[arrow] (dec)--node[above,font=\scriptsize]{بازسازی پیام}(dst);
  \node[box,below=0.68cm of ch,minimum width=1.5cm,minimum height=0.65cm,fill=warncolor!7,draw=warncolor!65] (noise) {منبع نویز};
  \draw[arrow,draw=warncolor] (noise.north)--(ch.south);
\end{tikzpicture}
\caption{مدل ساده‌ی ارتباط. نویز می‌تواند در مسیر کانال یا در بخش‌های دیگر سامانه وارد شود و سیگنال را تغییر دهد.}
\label{fig:communication-model}
\end{figure}

این مدل به ما کمک می‌کند سؤال‌های دقیق‌تری بپرسیم. آیا رمزگذار می‌تواند هر پیامی را به هر سیگنالی تبدیل کند؟ آیا کانال ظرفیت انتقال نامحدود دارد؟ اگر بخشی از سیگنال در طول مسیر تغییر کند، گیرنده چطور می‌تواند پیام اصلی را بازیابی کند؟ و آیا یک پیام را می‌توان به چند روش فیزیکی متفاوت فرستاد؟

پاسخ پرسش آخر مثبت است. همان جمله را می‌توان با جوهر روی کاغذ نوشت، با فشار هوا منتقل کرد، در قالب پالس‌های نوری فرستاد یا به شکل تغییرات ولتاژ در یک سیم نمایش داد. پس لازم نیست یک پیام همیشه یک حامل فیزیکی خاص داشته باشد. آنچه باید حفظ شود، رابطه‌ی قراردادی یا ساختاری میان پیام، کد و چیزی است که گیرنده دریافت می‌کند.

نکته‌ی دیگر این است که منبع پیام الزاماً یک انسان نیست. یک حسگر دما، یک دستگاه اندازه‌گیری، یک رایانه یا حتی یک فرایند طبیعی می‌تواند منبعی باشد که از میان چندین حالت ممکن، وضعیتی را تولید می‌کند. همین‌طور مقصد نیز الزاماً انسان نیست؛ ممکن است یک رایانه، یک کنترل‌کننده یا دستگاهی باشد که بر اساس سیگنال دریافتی تصمیمی می‌گیرد. به همین دلیل مدل شانون فقط برای مکالمه‌ی انسانی نیست و می‌تواند در بررسی سامانه‌های مخابراتی، ذخیره‌سازی و پردازش داده نیز به کار رود.

\subsection{معنای پیام با مقدار اطلاعات یکی نیست}
\label{sub:meaning-versus-information}

اینجا باید به نکته‌ای برگردیم که ممکن است در نگاه اول عجیب به نظر برسد. اگر جمله‌ای بسیار مهم درباره‌ی زندگی انسان دریافت کنیم، آیا مقدار اطلاعات آن از یک پیام ساده مثل «در باز است» بیشتر است؟ از دیدگاه انسانی ممکن است پاسخ مثبت باشد؛ اما نظریه‌ی اطلاعات شانون از این زاویه به پیام نگاه نمی‌کند.

شانون صریحاً توضیح می‌دهد که جنبه‌های معنایی ارتباط برای مسئله‌ی مهندسی انتقال، موضوع اصلی نیستند. آنچه در مدل او اهمیت دارد این است که پیام واقعی از میان مجموعه‌ای از پیام‌های ممکن انتخاب شده است. به بیان دیگر، نظریه‌ی او می‌پرسد پیام انتخاب‌شده چقدر از میان گزینه‌های ممکن قابل تمایز است و انتقال آن تا چه اندازه قابل اعتماد خواهد بود؛ نه اینکه آن پیام از نظر فلسفی، سیاسی یا شخصی چه معنایی دارد. \autocite{shannon1948}

برای مثال، فرض کنید یک فرستنده قرار است فقط یکی از دو نماد «۰» یا «۱» را بفرستد. اگر گیرنده نداند کدام نماد انتخاب شده، دریافت هر کدام از آن‌ها وضعیت او را تغییر می‌دهد. اما تصور کنید منبع همیشه و بدون استثنا نماد «۰» را می‌فرستد. در چنین سامانه‌ای، گیرنده حتی پیش از دریافت پیام هم می‌تواند حدس بزند که چه چیزی خواهد آمد. پس پیام از نظر معنا شاید برای کاربردی خاص مهم باشد، اما از نظر کاهش عدم‌قطعیت چیز تازه‌ای به گیرنده اضافه نمی‌کند.

این تفاوت را می‌توان با مثال آب‌وهوا هم دید. جمله‌ی «فردا خورشید طلوع می‌کند» برای بیشتر ما خبری غافلگیرکننده نیست. اما جمله‌ی «فردا در میانه‌ی روز خورشید برای چند ساعت ناپدید خواهد شد» ممکن است بسیار غیرمنتظره باشد. از نظر نظریه‌ی اطلاعات، غیرمنتظره‌بودن به احتمال رخداد مربوط است. بااین‌حال، این مسئله هنوز به‌معنای مهم‌تر یا معنادارتر بودن پیام دوم در تمام زمینه‌ها نیست؛ فقط می‌گوید اگر مدل احتمالاتی ما آن را بسیار نامحتمل بداند، وقوعش عدم‌قطعیت بیشتری را از میان می‌برد.

پس بهتر است از همین‌جا دو پرسش را از یکدیگر جدا کنیم: «این پیام برای ما چه معنایی دارد؟» و «این پیام در یک مدل احتمالاتی تا چه اندازه غیرمنتظره است؟» پرسش اول به زمینه، قرارداد، دانش و تفسیر مربوط می‌شود. پرسش دوم به احتمال و اندازه‌گیری اطلاعات در نظریه‌ی شانون. ممکن است این دو در کاربردهای واقعی با یکدیگر ارتباط پیدا کنند، اما یکسان نیستند.

\subsection{پیش از شانون: نایکوئیست و هارتلی}
\label{sub:nyquist-and-hartley}

نظریه‌ی اطلاعات یک‌باره و در خلأ به وجود نیامد. پیش از آنکه کلود شانون در سال ۱۹۴۸ صورت‌بندی مشهور خود را منتشر کند، مهندسان مخابرات با مسئله‌ای کاملاً عملی روبه‌رو بودند: چگونه می‌توان در یک بازه‌ی زمانی مشخص، بیشترین مقدار پیام را از یک خط تلگراف یا یک کانال ارتباطی عبور داد، بی‌آنکه سیگنال‌ها آن‌قدر در هم فرو بروند که گیرنده نتواند آن‌ها را از هم تشخیص دهد؟

هری نایکوئیست در مقاله‌ی سال ۱۹۲۴ خود، با عنوان \lr{Certain Factors Affecting Telegraph Speed}، دو مسئله را بررسی کرد: شکل‌دهی سیگنال و انتخاب کد. او نشان داد سرعت تلگراف فقط به این وابسته نیست که فرستنده چند بار بتواند وضعیت جریان را تغییر دهد؛ شکل موج، ویژگی‌های خط و میزان تداخل میان علامت‌ها نیز اهمیت دارند. همچنین تعداد حالت‌های متفاوتی که می‌توان برای ارسال علامت به کار برد، محدودیت‌هایی عملی دارد. بنابراین پیش از آنکه درباره‌ی «مقدار اطلاعات» صحبت کنیم، باید بفهمیم یک کانال فیزیکی اصلاً چه تعداد حالت قابل تشخیص را در اختیار ما می‌گذارد. \autocite{nyquist1924}

چهار سال بعد، رالف هارتلی در مقاله‌ی \lr{Transmission of Information} تلاش کرد مقدار اطلاعات را با کمیتی ریاضی بیان کند. ایده‌ی او این بود که اگر برای هر جایگاه از یک پیام، بتوانیم یکی از چند نماد را انتخاب کنیم، تعداد پیام‌های ممکن را می‌توانیم بشماریم. فرض کنید هر پیام از \(n\) جایگاه تشکیل شده باشد و در هر جایگاه یکی از \(s\) نماد مجاز قرار گیرد. در این صورت تعداد رشته‌های ممکن برابر است با:

\[
M=s^n
\]

اگر همه‌ی این پیام‌ها را هم‌احتمال فرض کنیم، مقدار اطلاعات لازم برای مشخص‌کردن یکی از آن‌ها برابر لگاریتم پایه‌ی دوِ تعداد پیام‌های ممکن است:

\[
I=\log_2 M=\log_2(s^n)=n\log_2 s
\]

برای مثال، اگر سه جایگاه داشته باشیم و در هر جایگاه یکی از چهار نماد را قرار دهیم، تعداد پیام‌های ممکن \(4^3=64\) است. چون \(\log_2 64=6\)، برای مشخص‌کردن یکی از این ۶۴ پیام به شش بیت اطلاعات نیاز داریم. در اینجا بیت هنوز فقط واحد ریاضی شمارش انتخاب‌هاست؛ نه سیمی در مدار و نه حالتی فیزیکی در حافظه. \autocite{hartley1928}

اما معیار هارتلی یک فرض مهم دارد: همه‌ی پیام‌های ممکن را هم‌احتمال در نظر می‌گیرد. اگر بعضی پیام‌ها بسیار بیشتر از پیام‌های دیگر رخ دهند، شمارش ساده‌ی آن‌ها نمی‌تواند عدم‌قطعیت واقعی منبع را به‌درستی توصیف کند. همین‌جا مسئله‌ای باقی می‌ماند که شانون به آن پرداخت: به‌جای اینکه فقط تعداد انتخاب‌ها را بشماریم، احتمال هر انتخاب را نیز وارد محاسبه کنیم. \autocite{shannon1948}

\section{عدم‌قطعیت: اطلاعات پیش از دریافت پیام}
\label{sec:uncertainty}

فرض کنید دوستی به شما می‌گوید: «یک عدد بین یک تا شش انتخاب کرده‌ام. حدس بزن کدام است.» اگر هیچ اطلاعات دیگری در اختیار نداشته باشید، شش انتخاب برایتان ممکن است. حالا فرض کنید او به‌جای این کار، یک سکه‌ی معمولی را پرتاب کرده و از شما می‌خواهد نتیجه را حدس بزنید. این بار فقط دو انتخاب دارید.

در نگاه اول شاید بگوییم مسئله‌ی اول اطلاعات بیشتری دارد، چون تعداد حالت‌های ممکن آن بیشتر است. اما اگر بدانیم دوستمان در تاس انداختن همیشه عدد شش می‌آورد، وضعیت تغییر می‌کند. تعداد حالت‌های ممکن هنوز شش تاست، ولی احتمال‌ها برابر نیستند. در واقع آنچه پیش از دریافت پیام می‌دانیم، فقط به تعداد گزینه‌ها وابسته نیست؛ به احتمال هر گزینه نیز بستگی دارد.

این‌جاست که احتمال وارد بحث می‌شود. احتمال زبانی برای توصیف انتظار ما از رخدادهای ممکن است. این زبان می‌تواند درباره‌ی یک آزمایش تکرارشونده، یک منبع پیام، یک اندازه‌گیری ناقص یا حتی دانسته‌های محدود ما درباره‌ی یک وضعیت فیزیکی صحبت کند. اما نباید بدون دقت، احتمال را با خودِ واقعیت یکی بگیریم.

\subsection{احتمال به‌مثابه‌ی زبان توصیف عدم‌قطعیت}
\label{sub:probability-as-language}

به یک سکه نگاه کنید. اگر سکه منصفانه باشد، احتمال شیر و خط هر کدام برابر با ۰٫۵ است. این یعنی اگر پرتاب را در شرایط مناسب بارها تکرار کنیم، انتظار داریم نسبت تعداد شیرها و خط‌ها به یکدیگر نزدیک شود. اما در یک پرتاب مشخص، نتیجه فقط یکی از آن دو خواهد بود.

حالا یک سکه‌ی نامتوازن را در نظر بگیرید که احتمال شیر در آن ۰٫۹ و احتمال خط ۰٫۱ است. باز هم دو نتیجه ممکن داریم، اما پیش‌بینی ما تغییر کرده است. احتمال شیر بیشتر است و اگر مجبور باشیم پیش از پرتاب یک نتیجه را انتخاب کنیم، انتخاب شیر منطقی‌تر خواهد بود. بااین‌حال، نتیجه‌ی خط همچنان ممکن است رخ دهد. احتمال پایین به‌معنای ناممکن‌بودن نیست.

در این مثال‌ها، احتمال می‌تواند دو نقش متفاوت داشته باشد. گاهی درباره‌ی فراوانی یک رخداد در مجموعه‌ای از آزمایش‌های تکرارشونده صحبت می‌کنیم؛ برای مثال، احتمال خط در یک سکه‌ی نامتوازن. گاهی هم احتمال را برای بیان میزان باور یا عدم‌قطعیت خود به کار می‌بریم؛ مثلاً وقتی هنوز نمی‌دانیم یک قطعه خراب است یا سالم. این دو برداشت در فلسفه‌ی احتمال و آمار تاریخ مفصلی دارند و همیشه هم بر سر تفسیرشان توافق وجود ندارد. برای بحث ما کافی است بدانیم احتمال بخشی از مدل توصیف وضعیت‌های ممکن است، نه برچسبی جادویی که روی خود اشیا چسبانده شده باشد. \autocite{jaynes2003,cover2006}

یک نکته‌ی مهم‌تر نیز وجود دارد. فرض کنید سکه در هوا در حال چرخش است. اگر تمام شرایط اولیه، سرعت، زاویه، مقاومت هوا و ویژگی‌های مکانیکی را با دقتی نامحدود می‌دانستیم، شاید می‌توانستیم نتیجه را در چارچوب مدل کلاسیکی پیش‌بینی کنیم. اما در عمل چنین دانشی نداریم. بنابراین احتمال می‌تواند عدم‌قطعیت ناشی از محدودیت اطلاعات ما را توصیف کند، بی‌آنکه الزاماً ادعا کنیم خودِ رویداد از نظر فیزیکی نامعین است. در مکانیک کوانتومی مسئله شکل متفاوت و عمیق‌تری پیدا می‌کند، اما برای ورود به نظریه‌ی اطلاعات لازم نیست تمام این بحث را یک‌جا حل کنیم.

پس وقتی در ادامه از احتمال یک پیام حرف می‌زنیم، باید بدانیم که احتمال را نسبت به چه مدلی تعریف کرده‌ایم. اگر منبعی را بررسی می‌کنیم که گاهی «۰» و گاهی «۱» تولید می‌کند، باید توزیع احتمال این دو خروجی را مشخص کنیم. اگر احتمال‌ها را ندانیم یا اشتباه برآورد کنیم، اندازه‌گیری اطلاعات نیز به همان نسبت تغییر می‌کند.

\subsection{از غافلگیری تا آنتروپی شانون}
\label{sub:shannon-entropy}

حالا به پرسش اصلی برگردیم: چگونه می‌توانیم مقدار اطلاعات حاصل از یک رخداد را با عددی مشخص بیان کنیم؟ یک راه شهودی این است که بگوییم هرچه یک رخداد نامحتمل‌تر باشد، وقوع آن اطلاعات بیشتری به ما می‌دهد. اگر بدانیم خورشید هر روز طلوع می‌کند و یک روز هم طلوع کند، تعجب نمی‌کنیم. اما اگر رخدادی بسیار نامحتمل اتفاق بیفتد، احتمالاً انتظار ما به‌شدت تغییر می‌کند.

برای بیان این ایده می‌توانیم مقدار اطلاعات یک رخداد با احتمال \(p(x)\) را به شکل زیر تعریف کنیم:

\[
I(x)=-\log_2 p(x)
\]

این کمیت را خوداطلاعاتی یا اطلاعاتِ رخداد می‌نامند. اگر احتمال رخداد برابر با یک باشد، آنگاه \(I(x)=0\) است؛ چون رخداد کاملاً قابل پیش‌بینی بوده و وقوع آن چیزی از عدم‌قطعیت ما کم نمی‌کند. اگر احتمال رخداد نصف باشد، مقدار آن یک بیت خواهد بود. و هرچه احتمال کمتر شود، مقدار خوداطلاعاتی افزایش می‌یابد.

چرا لگاریتم؟ یکی از دلایل، همان خاصیتی است که پیش‌تر در کار هارتلی هم دیدیم. اگر دو رخداد مستقل داشته باشیم، احتمال وقوع هم‌زمان آن‌ها برابر حاصل‌ضرب احتمال‌هایشان است. لگاریتم این ضرب را به جمع تبدیل می‌کند. اگر دو انتخاب مستقل هر کدام یک بیت اطلاعات داشته باشند، مجموع آن‌ها دو بیت می‌شود؛ نه اینکه مجبور باشیم تعداد حالت‌های ترکیبی را جداگانه بشماریم. این ویژگی برای اندازه‌گیری رشته‌های طولانی از نمادها بسیار مفید است. \autocite{hartley1928,shannon1948}

اما یک رخداد مشخص با کل منبع پیام فرق دارد. پیش از اینکه پیام را دریافت کنیم، ممکن است چند نتیجه‌ی مختلف با احتمال‌های متفاوت در انتظار ما باشند. بنابراین به‌جای اینکه فقط اطلاعات یک نتیجه را اندازه بگیریم، می‌خواهیم مقدار متوسط اطلاعاتی را که از خروجی‌های منبع به دست می‌آید محاسبه کنیم. شانون این کمیت را با آنتروپی نشان داد:

\[
H(X)=-\sum_{i=1}^{n}p(x_i)\log_2 p(x_i)
\]

در این رابطه، \(X\) متغیری است که می‌تواند یکی از چند نتیجه‌ی ممکن را بگیرد و \(p(x_i)\) احتمال نتیجه‌ی \(x_i\) است. آنتروپی، میانگین خوداطلاعاتی نتیجه‌هایی است که منبع می‌تواند تولید کند. وقتی لگاریتم پایه‌ی دو باشد، واحد اندازه‌گیری بیت است. \autocite{shannon1948,cover2006}

برای یک سکه‌ی منصفانه داریم:

\[
H(X)=-\left(0.5\log_2 0.5+0.5\log_2 0.5\right)=1\ \text{\lr{bit}}
\]

اما برای سکه‌ای که احتمال شیر در آن ۰٫۹ و احتمال خط ۰٫۱ است:

\[
H(X)=-\left(0.9\log_2 0.9+0.1\log_2 0.1\right)\approx0.469\ \text{\lr{bit}}
\]

در هر دو سکه دو نتیجه‌ی ممکن داریم، اما آنتروپی آن‌ها یکسان نیست. در سکه‌ی منصفانه، پیش از پرتاب عدم‌قطعیت بیشتری داریم؛ چون هیچ‌کدام از دو نتیجه بر دیگری برتری احتمالی ندارد. در سکه‌ی نامتوازن، شیر بسیار محتمل‌تر است و به همین دلیل میانگین اطلاعات حاصل از هر پرتاب کمتر می‌شود.

حالا به‌جای سکه، یک تاس چهاروجهی را در نظر بگیریم. اگر تاس منصفانه باشد، هر چهار وجه احتمال یکسانی دارند؛ یعنی احتمال هر نتیجه یک‌چهارم است. در این حالت آنتروپی برابر می‌شود با:

\[
H(X)=-4\left(\frac14\log_2\frac14\right)=2\;\text{bit}
\]

اما فرض کنید تاس دیگری داریم که چهار وجه آن احتمال‌های یکسانی ندارند: احتمال وجه اول یک‌دوم، وجه دوم یک‌چهارم، و احتمال دو وجه دیگر هرکدام یک‌هشتم است. آنتروپی این تاس چنین خواهد بود:

\[
\begin{aligned}
H(X)={}&-\frac12\log_2\frac12
        -\frac14\log_2\frac14\\
      &-\frac18\log_2\frac18
        -\frac18\log_2\frac18\\
    ={}&0.5+0.5+0.375+0.375\\
    ={}&1.75\;\text{bit}
\end{aligned}
\]

هر دو تاس چهار نتیجه‌ی ممکن دارند، اما آنتروپی تاس منصفانه دو بیت و آنتروپی تاس نامتوازن ۱٫۷۵ بیت است. تفاوت از تعداد وجه‌ها نمی‌آید؛ از توزیع احتمال آن‌ها می‌آید. وقتی احتمال‌ها برابر باشند، پیش از دیدن نتیجه بیشترین عدم‌قطعیت را داریم. هرچه احتمال‌ها نامتوازن‌تر شوند، پیش‌بینی نتیجه آسان‌تر می‌شود و آنتروپی کاهش می‌یابد. \autocite{cover2006}

\begin{figure}[htbp]
\centering
\begin{tikzpicture}[x=0.95cm,y=0.72cm]
  \node[anchor=east] at (-0.15,2.8) {سکه‌ی منصفانه};
  \fill[genixblue!75] (0,2.2) rectangle (2,2.8);
  \fill[genixteal!75] (2,2.2) rectangle (4,2.8);
  \node[white] at (1,2.5) {شیر ۰٫۵};
  \node[white] at (3,2.5) {خط ۰٫۵};
  \node[anchor=west] at (4.2,2.5) {آنتروپی: ۱ بیت};
  \node[anchor=east] at (-0.15,1.1) {سکه‌ی نامتوازن};
  \fill[genixblue!75] (0,0.5) rectangle (3.6,1.1);
  \fill[genixteal!75] (3.6,0.5) rectangle (4,1.1);
  \node[white] at (1.8,0.8) {شیر ۰٫۹};
  \node[black] at (3.8,0.8) {۰٫۱};
  \node[anchor=west] at (4.2,0.8) {آنتروپی: ۰٫۴۶۹ بیت};
\end{tikzpicture}
\caption{تعداد نتیجه‌های ممکن در دو سکه برابر است، اما توزیع احتمال آن‌ها متفاوت است؛ بنابراین آنتروپی یکسانی ندارند.}
\label{fig:coin-entropy}
\end{figure}

از اینجا می‌توانیم چند نکته را بیرون بکشیم. نخست اینکه آنتروپی به توزیع احتمال وابسته است، نه فقط تعداد نتیجه‌های ممکن. دوم اینکه آنتروپی درباره‌ی متوسط عدم‌قطعیت منبع است، نه درباره‌ی معنای پیام. سوم اینکه آنتروپی بالا لزوماً نشانه‌ی مفیدتر، عمیق‌تر یا معنادارتر بودن داده نیست. یک رشته‌ی تصادفی می‌تواند آنتروپی آماری بالایی داشته باشد، اما از نظر انسانی هیچ پیام خاصی برای ما نداشته باشد.

\subsection{آنتروپی با اطلاعات یک پیام چه تفاوتی دارد؟}
\label{sub:entropy-and-message-information}

ممکن است تا اینجا چنین به نظر برسد که اگر آنتروپی یک منبع یک بیت باشد، هر پیامی که از آن دریافت می‌کنیم دقیقاً یک بیت اطلاعات دارد. اما این برداشت درست نیست. آنتروپی میانگین است. یک پیام مشخص می‌تواند از نظر خوداطلاعاتی بیشتر یا کمتر از میانگین باشد.

برای مثال، در سکه‌ی نامتوازن، رخداد شیر بسیار محتمل است و خوداطلاعاتی کمتری دارد. رخداد خط نامحتمل‌تر است و خوداطلاعاتی بیشتری دارد. اما آنتروپی، میانگین این دو حالت با در نظر گرفتن احتمال رخدادشان است. به همین دلیل است که گاهی یک پیام خاص بسیار غافلگیرکننده است، درحالی‌که منبعی که آن را تولید کرده به‌طور متوسط آنتروپی پایینی دارد.

در نظریه‌ی اطلاعات کمیت دیگری نیز داریم که برای ارتباط مهم است: اطلاعات متقابل. فرض کنید منبعی نماد \(X\) را تولید می‌کند و گیرنده پس از عبور پیام از کانال، خروجی \(Y\) را دریافت می‌کند. اطلاعات متقابل به ما می‌گوید دانستن خروجی کانال تا چه اندازه عدم‌قطعیت ما درباره‌ی ورودی را کاهش می‌دهد:

\[
I(X;Y)=H(X)-H(X|Y)
\]

در این رابطه، \(H(X|Y)\) عدم‌قطعیت باقی‌مانده درباره‌ی ورودی پس از مشاهده‌ی خروجی است. اگر خروجی هیچ ارتباطی با ورودی نداشته باشد، دانستن آن کمکی به ما نمی‌کند و اطلاعات متقابل صفر می‌شود. اگر خروجی بتواند ورودی را با اطمینان کامل مشخص کند، اطلاعات متقابل می‌تواند به آنتروپی ورودی برسد. \autocite{cover2006}

این مفهوم برای فهم ارتباط بسیار مهم است، چون مسئله فقط این نیست که منبع چقدر اطلاعات تولید می‌کند؛ مسئله این هم هست که چه مقدار از آن اطلاعات واقعاً به مقصد می‌رسد. اگر کانال پر از اختلال باشد، ممکن است منبع پیام‌های متنوعی تولید کند، اما گیرنده نتواند تفاوت میان آن‌ها را تشخیص دهد. در این حالت، مقدار اطلاعاتی که از دید گیرنده منتقل شده کمتر از چیزی است که منبع تولید کرده است.

\section{کانال، نویز و ظرفیت: چه مقدار از پیام واقعاً می‌رسد؟}
\label{sec:channel-noise-capacity}

تا اینجا فرض کردیم که پیام از منبع به گیرنده می‌رسد و گیرنده می‌تواند نمادهای دریافت‌شده را بخواند. اما در جهان واقعی، ارتباط به این سادگی نیست. وقتی با دو نفر در یک اتاق صحبت می‌کنیم، ممکن است صدای کولر، رفت‌وآمد افراد، فاصله، یا حتی صدای خودمان مانع شنیدن درست جمله شود. وقتی یک پیام را از طریق تلفن می‌فرستیم، ممکن است بخشی از صدا قطع شود. وقتی داده را از یک کابل عبور می‌دهیم، ویژگی‌های خود کابل و محیط پیرامون آن می‌توانند شکل سیگنال را تغییر دهند.

پس باید یک عنصر دیگر را وارد مدل کنیم: \textbf{کانال}. کانال همان مسیری است که پیام از طریق آن از فرستنده به گیرنده منتقل می‌شود. این مسیر می‌تواند هوا، یک رشته سیم، فیبر نوری، یک موج رادیویی، یا حتی یک سامانه‌ی پیچیده‌ی ارتباطی باشد. اما کانال فقط یک مسیر خالی نیست؛ ویژگی‌های آن تعیین می‌کنند چه نوع سیگنالی را می‌توانیم از آن عبور دهیم، با چه سرعتی، و با چه میزان خطا.

در نظریه‌ی اطلاعات، کانال را معمولاً با رابطه‌ای احتمالاتی توصیف می‌کنیم: اگر ورودی کانال را بدانیم، احتمال مشاهده‌ی هر خروجی چقدر است؟ به زبان ساده، کانال به ما می‌گوید یک نماد ورودی ممکن است با چه احتمالی به یک نماد خروجی تبدیل شود. اگر کانال کاملاً قابل اعتماد باشد، هر نماد ورودی همیشه همان‌طور که فرستاده شده دریافت می‌شود. اما در یک کانال نویزی، یک نماد ممکن است اشتباه دریافت شود، حذف شود، یا با نمادی دیگر اشتباه گرفته شود. \autocite{shannon1948,cover2006}

\subsection{نویز دقیقاً چیست؟}
\label{sub:what-is-noise}

بیایید از یک مثال روزمره شروع کنیم. فرض کنید در یک ایستگاه قطار ایستاده‌اید و دوستتان از چند متر آن‌طرف‌تر نام شما را صدا می‌زند. در همان لحظه، قطاری وارد ایستگاه می‌شود. صدای قطار به صدای دوستتان اضافه می‌شود و ممکن است شما نام خود را اشتباه بشنوید. حالا سؤال این است: آیا صدای قطار بخشی از پیام دوست شماست؟ خیر. پیام همان چیزی است که دوستتان قصد داشته منتقل کند؛ صدای قطار اختلالی است که دریافت آن پیام را دشوار کرده است.

در مهندسی ارتباطات، به هر مؤلفه‌ی ناخواسته‌ای که دریافت یا تشخیص سیگنال را دشوار کند، به‌طور کلی نویز می‌گوییم. البته نویز همیشه صدایی نیست که با گوش بشنویم. در یک مدار الکترونیکی، تغییرات ناخواسته‌ی ولتاژ می‌تواند نقش نویز را داشته باشد. در ارتباط رادیویی، سیگنال‌های دیگر یا آشفتگی‌های محیطی می‌توانند در دریافت پیام اختلال ایجاد کنند. در یک تصویر دیجیتال، نقاط رنگی ناخواسته ممکن است جزئیاتی را بپوشانند که قرار بوده ثبت شوند.

نکته‌ی مهم این است که نویز را نباید با «اطلاعاتِ نامطلوب» یا «پیامی که دوست نداریم» یکی بدانیم. نویز نسبت به هدف و سامانه‌ی ارتباطی تعریف می‌شود. صدای قطار برای گفت‌وگوی دو نفر نویز است، اما برای کسی که در حال بررسی صدای قطار است، همان صدا خودِ داده‌ی موردنظر اوست. یک تغییر ولتاژ نیز ممکن است برای یک گیرنده اختلال باشد، اما در دستگاهی دیگر همان تغییر، پیام اصلی را حمل کند.

به همین دلیل، تشخیص نویز فقط به خودِ پدیده وابسته نیست؛ باید بدانیم چه چیزی قرار بوده منتقل شود و گیرنده چگونه آن را تفسیر می‌کند. در واقع، یک رخداد فیزیکی به‌خودی‌خود برچسب «پیام» یا «نویز» ندارد. این برچسب را نسبت آن رخداد با سامانه‌ی ارتباطی مشخص می‌کند.

\subsection{چرا نمی‌توان هرقدر خواست اطلاعات فرستاد؟}
\label{sub:channel-capacity}

حالا فرض کنید می‌خواهیم از همان ایستگاه قطار پیامی را منتقل کنیم. اگر نام دوستمان را آرام و شمرده بگوییم، احتمالاً او آن را می‌شنود. اما اگر بخواهیم در یک ثانیه صدها واژه را فریاد بزنیم، حتی بدون نویز هم ممکن است نتواند آن‌ها را از هم تشخیص دهد. پس محدودیت فقط از نویز نمی‌آید؛ خودِ کانال هم ظرفیت محدودی دارد.

ظرفیت کانال، به‌طور خلاصه، بیشترین نرخ اطلاعاتی است که می‌توان با نرخ خطای دلخواهِ کوچک از یک کانال منتقل کرد. این تعریف به جزئیات ریاضی بیشتری نیاز دارد، اما فعلاً ایده‌ی اصلی کافی است: هر کانال در شرایط مشخص، سقفی برای انتقال قابل اعتماد اطلاعات دارد. اگر بخواهیم سریع‌تر از آن سقف پیام بفرستیم، باید هزینه‌ی آن را به شکل خطای بیشتر، نیاز به کدگذاری پیچیده‌تر، یا منابع فیزیکی بیشتر بپذیریم. \autocite{shannon1948,cover2006}

برای روشن‌تر شدن موضوع، یک مثال ساده را در نظر بگیریم. فرض کنید کانالی داریم که در هر لحظه فقط می‌تواند یکی از دو نماد صفر یا یک را منتقل کند. اگر این کانال در هر ثانیه هزار بار بتواند یک نماد را منتقل کند، نرخ خام انتقال آن هزار نماد در ثانیه است. اما اگر کانال نویز زیادی داشته باشد و گیرنده مرتب صفر را به‌جای یک یا یک را به‌جای صفر دریافت کند، نرخ انتقال خام به‌تنهایی چیزی درباره‌ی ارتباط قابل اعتماد به ما نمی‌گوید. ممکن است هزار نماد در ثانیه بفرستیم، ولی بخش قابل توجهی از پیام به مقصد نرسد.

اینجا باید میان سه چیز تفاوت بگذاریم: تعداد نمادهایی که در واحد زمان فرستاده می‌شوند، مقدار اطلاعاتی که منبع در هر نماد تولید می‌کند، و مقدار اطلاعاتی که گیرنده در نهایت می‌تواند با اطمینان بازیابی کند. این سه کمیت همیشه یکسان نیستند.

برای یک نمونه‌ی مشخص از کانال، یعنی کانال با پهنای باند محدود و نویز سفید گوسی افزایشی، رابطه‌ای مشهور به نام قضیه‌ی شانون و هارتلی داریم:

\[
C=B\log_2\left(1+\frac{S}{N}\right)
\]

در این رابطه، \(C\) ظرفیت نظری کانال بر حسب بیت بر ثانیه، \(B\) پهنای باند بر حسب هرتز، و \(S/N\) نسبت توان سیگنال به توان نویز است. این فرمول نشان می‌دهد ظرفیت انتقال فقط به سرعت تغییر سیگنال وابسته نیست؛ پهنای باند و نسبت سیگنال به نویز نیز اهمیت دارند. اگر پهنای باند بیشتر شود، یا سیگنال در مقایسه با نویز قوی‌تر باشد، ظرفیت نظری کانال افزایش پیدا می‌کند. \autocite{cover2006,shannon1948}

اما حواسمان باشد: این رابطه قانون همگانی تمام کانال‌های ارتباطی نیست. این فرمول برای مدل مشخصی از کانال، با فرض‌های معین، به دست آمده است. کانال‌های واقعی ممکن است محدودیت‌ها و ویژگی‌های دیگری داشته باشند؛ برای همین نباید آن را بدون توجه به مدل، به هر کابل، موج رادیویی یا سامانه‌ی ارتباطی تعمیم دهیم.

\subsection{کدگذاری: چرا گاهی پیام را طولانی‌تر می‌کنیم؟}
\label{sub:redundancy-and-error-correction}

در نگاه اول عجیب است که برای انتقال بهتر یک پیام، گاهی باید آن را طولانی‌تر کنیم. فرض کنید می‌خواهید یک بیت را از یک کانال نویزی عبور دهید. اگر فقط یک بار آن را بفرستید، ممکن است همان یک ارسال اشتباه دریافت شود. حالا به‌جای یک بار، سه بار همان بیت را ارسال می‌کنید: برای صفر، «۰۰۰» و برای یک، «۱۱۱». گیرنده نیز با رأی اکثریت تصمیم می‌گیرد که پیام اصلی چه بوده است.

برای مثال اگر «۱۱۱» ارسال شده باشد و یکی از بیت‌ها به‌اشتباه صفر دریافت شود، خروجی «۱۰۱» هنوز با رأی اکثریت به یک تعبیر می‌شود. این روش ساده، نرخ انتقال اطلاعات را کاهش می‌دهد، چون برای یک بیت اطلاعاتی سه نماد فرستاده‌ایم؛ اما در عوض احتمال خطای نهایی را کمتر می‌کند.

این همان ایده‌ی کلی افزونگی است: بخشی از آنچه ارسال می‌کنیم، تکرار یا ساختاری است که به گیرنده کمک می‌کند خطا را تشخیص دهد یا اصلاح کند. در زبان عادی، شاید تکرار را بیهوده بدانیم؛ اما در ارتباطات، تکرار حساب‌شده می‌تواند تفاوت میان یک پیام قابل بازیابی و یک پیام خراب‌شده باشد. کدهای تصحیح خطا این ایده را به‌شکل ریاضی و بسیار کارآمدتری پیاده می‌کنند.

در اینجا یک ظرافت وجود دارد. کدگذاری نمی‌تواند ظرفیت فیزیکی کانال را جادویی حذف کند. کاری که انجام می‌دهد این است که از ظرفیت موجود هوشمندانه استفاده می‌کند و با افزودن ساختار مناسب، احتمال خطا را کاهش می‌دهد. قضیه‌ی بنیادی شانون درباره‌ی کانال‌های نویزی دقیقاً نشان می‌دهد که برای نرخ‌هایی پایین‌تر از ظرفیت، می‌توان با کدگذاری مناسب احتمال خطا را هرقدر بخواهیم کوچک کرد؛ اما بالاتر از ظرفیت، چنین تضمینی ممکن نیست. \autocite{shannon1948,cover2006}

\subsection{آنتروپی شانون و پیچیدگی کولموگروف: دو پرسش متفاوت}
\label{sub:kolmogorov-complexity}

تا اینجا بارها از آنتروپی شانون صحبت کردیم. اما مفهوم دیگری هم در ریاضیات و علوم رایانه وجود دارد که نامش به اطلاعات گره خورده است: پیچیدگی کولموگروف. شاید در نگاه اول این دو مفهوم بسیار شبیه به نظر برسند، چون هر دو به نوعی درباره‌ی مقدار اطلاعات صحبت می‌کنند. بااین‌حال، آن‌ها یک پرسش را پاسخ نمی‌دهند.

آنتروپی شانون به یک منبع احتمالاتی نگاه می‌کند و می‌پرسد: اگر خروجی‌های این منبع را بارها مشاهده کنیم، به‌طور متوسط با چه میزان عدم‌قطعیتی روبه‌رو خواهیم بود؟ اما پیچیدگی کولموگروف به یک رشته‌ی مشخص نگاه می‌کند و می‌پرسد: کوتاه‌ترین برنامه‌ای که بتواند دقیقاً همین رشته را تولید کند، چقدر طول دارد؟ \autocite{kolmogorov1965,solomonoff1964,chaitin1966,li-vitanyi2008}

برای مثال دو رشته را در نظر بگیرید:

\[
A=01010101010101010101
\]

\[
B=01101000110101110010
\]

هر دو رشته بیست بیت طول دارند. اگر فقط خود رشته‌ها را ببینیم، تعداد بیت‌هایشان برابر است. اما برای توصیف رشته‌ی اول می‌توانیم یک برنامه‌ی بسیار کوتاه بنویسیم: «۱۰ را ده بار تکرار کن». در مقابل، اگر رشته‌ی دوم هیچ الگوی ساده‌ای نداشته باشد، ممکن است برای تولید دقیق آن ناچار شویم تقریباً تمام بیت‌هایش را در برنامه بیاوریم. پس رشته‌ی اول از نظر توصیف الگوریتمی ساده‌تر است، درحالی‌که رشته‌ی دوم ممکن است پیچیدگی بیشتری داشته باشد.

به این معنا، پیچیدگی کولموگروف به فشرده‌پذیری الگوریتمی یک داده مربوط است. اگر بتوانیم داده‌ای را با یک دستور کوتاه تولید کنیم، توصیف کوتاه‌تری برای آن داریم. اگر هیچ راه کوتاه‌تری جز نوشتن خود داده وجود نداشته باشد، توصیف آن تقریباً به اندازه‌ی خود داده طول خواهد داشت. البته تعریف دقیق این مفهوم به ماشین محاسباتی عمومی و زبان برنامه‌نویسی فرضی وابسته است و مقدار آن تا یک ثابت مستقل از طول رشته تغییر می‌کند. \autocite{li-vitanyi2008}

اما نباید از این توضیح نتیجه بگیریم که رشته‌ای که ظاهراً تصادفی به نظر می‌رسد حتماً پیچیدگی کولموگروف بالایی دارد. ممکن است رشته‌ای طولانی و در نگاه اول نامنظم باشد، اما با یک قاعده‌ی کوتاه ساخته شده باشد؛ یا برعکس، الگویی در آن ببینیم که صرفاً تصادفی به نظر نمی‌رسد. علاوه‌براین، پیچیدگی کولموگروف برای رشته‌های دلخواه به‌طور عمومی قابل محاسبه نیست. می‌توانیم کران‌هایی برای آن پیدا کنیم یا درباره‌ی نمونه‌های خاص استدلال کنیم، اما الگوریتمی عمومی نداریم که برای هر رشته، کوتاه‌ترین برنامه‌ی ممکن را تضمین‌شده پیدا کند. \autocite{li-vitanyi2008}

پس آنتروپی شانون و پیچیدگی کولموگروف را نباید دو نام برای یک چیز بدانیم. اولی یک ویژگی آماریِ منبع یا توزیع احتمال است؛ دومی ویژگی الگوریتمیِ یک رشته‌ی مشخص. ممکن است یک منبع آنتروپی بالایی داشته باشد، اما رشته‌ای که در یک نوبت تولید کرده به‌دلیل وجود الگویی کوتاه، پیچیدگی الگوریتمی کمی داشته باشد. همچنین ممکن است یک رشته‌ی مشخص پیچیده باشد، بدون آنکه از همین موضوع بتوان درباره‌ی معنای آن نتیجه گرفت.

این تمایز برای ادامه‌ی بحث ما اهمیت دارد. وقتی از اطلاعات حرف می‌زنیم، باید بدانیم درباره‌ی کدام پرسش صحبت می‌کنیم: عدم‌قطعیت یک منبع، مقدار اطلاعات منتقل‌شده از یک کانال، یا کوتاهی و بلندی توصیف الگوریتمی یک داده؟ اگر این پرسش‌ها را با هم مخلوط کنیم، واژه‌ی «اطلاعات» به مفهومی مبهم تبدیل می‌شود؛ درست همان ابهامی که در ابتدای کتاب تلاش کردیم از آن فاصله بگیریم.

\section{بازنمایی فیزیکی چگونه ممکن است؟}
\label{sec:physical-representation}

در فصل قبل به این نتیجه رسیدیم که یک بازنمایی نمی‌تواند تا ابد در سطحی کاملاً انتزاعی باقی بماند. یک کلمه، یک عدد یا یک تصویر ممکن است در ذهن ما معنایی داشته باشد، اما اگر قرار باشد آن را به شخص دیگری منتقل کنیم، در جایی ثبت کنیم، یا بعداً دوباره بازیابی کنیم، باید چیزی وجود داشته باشد که این بازنمایی را در جهان قابل دسترس کند.

بیایید یک مثال ساده را در نظر بگیریم. فرض کنید می‌خواهم عدد ۵ را برای شما بفرستم. می‌توانم آن را روی کاغذ بنویسم، با صدای بلند بگویم، روی تخته نقاشی کنم، با انگشتانم نشان دهم، یا حتی با چراغی که پنج بار روشن و خاموش می‌شود نمایش دهم. در تمام این حالت‌ها، عددی که قصد انتقالش را دارم یکسان است، اما چیزی که آن را حمل می‌کند تغییر می‌کند.

در یک حالت با جوهر و کاغذ روبه‌رو هستیم؛ در حالت دیگر با موج صوتی؛ در حالت سوم با نور؛ و در حالت چهارم با حرکت دست. پس عدد ۵ همان جوهر، صدا، نور یا حرکت نیست. این‌ها راه‌های متفاوتی‌اند که به‌واسطه‌ی آن‌ها می‌توانیم یک محتوای مشخص را در جهان فیزیکی بازنمایی کنیم.

این مثال ظاهراً ساده، ما را به مسئله‌ای عمیق‌تر می‌رساند: چگونه ممکن است یک محتوای واحد در چند وضعیت فیزیکی کاملاً متفاوت تحقق پیدا کند؟ و از کجا می‌فهمیم که یک وضعیت فیزیکی مشخص، واقعاً نماینده‌ی همان عدد یا پیام است؟

\subsection{از نماد به حالت}
\label{sub:symbol-to-state}

پیش از هر چیز باید میان «نماد» و «چیزی که نماد را در جهان ثبت می‌کند» تفاوت بگذاریم. وقتی عدد ۵ را روی کاغذ می‌نویسیم، آنچه روی کاغذ وجود دارد مجموعه‌ای از خطوط و تغییرات رنگ است. اگر با میکروسکوپ به آن نگاه کنیم، دیگر با یک عدد مرتب و واضح روبه‌رو نیستیم؛ با رشته‌هایی از مولکول‌های جوهر، الیاف کاغذ و تغییرات میکروسکوپی سطح آن مواجه می‌شویم.

اما چرا ما این وضعیت فیزیکی را عدد ۵ می‌خوانیم؟ چون در یک نظام نمادگذاری، شکل مشخصی از خطوط را به مفهوم یا عدد مشخصی مربوط کرده‌ایم. اگر همان شکل را روی صفحه‌ی نمایش ببینیم، وضعیت فیزیکی کاملاً متفاوت است: نور از نقاط مختلف صفحه به چشم ما می‌رسد، و این الگوی نور در چارچوب قراردادها و سازوکار خواندن ما، به‌عنوان نماد ۵ تشخیص داده می‌شود.

در اینجا دو لایه را باید از هم جدا کنیم. از یک طرف، یک وضعیت فیزیکی مشخص داریم: جوهر روی کاغذ یا نوری که از صفحه می‌آید. از طرف دیگر، یک نظام تفسیر داریم که به آن وضعیت، نقش نمادین می‌دهد. خودِ خطوط جوهر هیچ ویژگی طبیعی و ذاتی‌ای ندارند که مجبورمان کند آن‌ها را عدد پنج بخوانیم. در زبان دیگری ممکن است شکل مشابه معنای دیگری داشته باشد، و در یک خط دیگر اصلاً نماد محسوب نشود.

پس یک وضعیت فیزیکی، تنها به‌دلیل شکل ظاهری‌اش بازنمایی نمی‌شود. باید رابطه‌ای میان آن وضعیت، چیزی که نمایندگی می‌کند، و سامانه‌ای که آن را تشخیص یا تفسیر می‌کند وجود داشته باشد. این همان نقطه‌ای است که بحث بازنمایی از سطح فلسفیِ صرف، به مسئله‌ی نشانه‌ها و ارتباط نزدیک می‌شود.

این موضوع را می‌توان با یک عبارت ساده خلاصه کرد: \textbf{نماد چیزی نیست که صرفاً وجود داشته باشد؛ نماد چیزی است که در یک رابطه، نقش نمایندگی پیدا می‌کند.} یک خط جوهر می‌تواند لکه باشد، اثر تصادفی قلم باشد، یا بخشی از یک حرف؛ اینکه کدام است به زمینه، قرارداد و تفسیر وابسته است.

\subsection{سه‌گانه‌ی پیرس: شمایل، نمایه و نماد}
\label{sub:peirce-triad}

برای روشن‌تر شدن رابطه‌ی میان نشانه و چیزی که به آن اشاره می‌کند، می‌توانیم به چارلز سندرس پیرس، فیلسوف و منطق‌دان آمریکایی، بازگردیم. پیرس نظریه‌ای درباره‌ی نشانه‌ها ارائه کرد که در آن نشانه صرفاً یک شیء یا شکل نیست؛ بلکه در رابطه‌ای سه‌گانه معنا پیدا می‌کند: نشانه، موضوعی که نشانه به آن مربوط است، و تفسیر یا اثری که این رابطه در ذهن یا فرایند تفسیرکننده ایجاد می‌کند. او این سه عنصر را به‌ترتیب با مفاهیم \lr{\textit{Sign}}، \lr{\textit{Object}} و \lr{\textit{Interpretant}} توضیح می‌دهد. \autocite{peirce-cp,sep-peirce-semiotics}

ممکن است بپرسید چرا پیرس به‌جای یک رابطه‌ی دوطرفه، بر سه عنصر تأکید می‌کند؟ چون صرف اینکه چیزی به چیز دیگری مربوط باشد، هنوز توضیح نمی‌دهد که این رابطه چگونه فهمیده می‌شود. برای اینکه یک نشانه واقعاً در یک فرایند دلالت عمل کند، باید تفسیر یا اثری نیز در کار باشد. برای مثال، دود را در نظر بگیرید. دود می‌تواند نشانه‌ی وجود آتش باشد؛ اما این رابطه زمانی برای ما معنا پیدا می‌کند که دود را تشخیص دهیم و آن را به احتمال وجود آتش مربوط کنیم.

پیرس برای طبقه‌بندی نشانه‌ها سه نوع مشهور را از هم جدا می‌کند: شمایل، نمایه و نماد. این سه دسته به ما کمک می‌کنند بفهمیم نشانه‌ها از چه راهی به موضوع خود مربوط می‌شوند.

\subsubsection{شمایل: نشانه از راه شباهت}

شمایل یا \lr{\textit{Icon}} نشانه‌ای است که به‌واسطه‌ی شباهت یا همانندی با موضوع خود عمل می‌کند. برای مثال، تصویر یک چهره می‌تواند شمایلی از آن شخص باشد؛ نقشه‌ی ساده‌ی یک اتاق می‌تواند از نظر آرایش فضایی، شمایلی از همان اتاق باشد؛ و تصویر یک درخت می‌تواند به‌خاطر شباهت ظاهری، ما را به یاد درخت واقعی بیندازد.

البته شباهت در اینجا لزوماً به معنای یکسان بودن نیست. یک نقشه‌ی مترو رنگ‌ها و خطوطی دارد که با تونل‌ها و قطارهای واقعی شباهت فیزیکی چندانی ندارند، اما ممکن است در روابط مکانی و مسیرها، ساختاری مشابه ایجاد کنند. بنابراین یک شمایل می‌تواند فقط در بعضی ویژگی‌ها شبیه موضوعش باشد، نه در تمام جزئیات.

\subsubsection{نمایه: نشانه از راه پیوند واقعی}

نمایه یا \lr{\textit{Index}} نشانه‌ای است که از طریق نوعی ارتباط واقعی، معمولاً علّی یا مکانی، با موضوع خود پیوند دارد. دود نشانه‌ی آتش است، ردپا نشانه‌ی عبور یک موجود است، و تب می‌تواند نشانه‌ای از تغییری در وضعیت بدن باشد. این نشانه‌ها لزوماً شبیه موضوع خود نیستند؛ دود شبیه آتش نیست و ردپا شبیه پای کسی نیست. رابطه‌ی آن‌ها از راه پیوندی دیگر برقرار می‌شود.

در این دسته، رابطه‌ی نشانه و موضوع اغلب به رخدادی در جهان وابسته است. اگر هیچ آتشی وجود نداشته باشد، دود واقعیِ ناشی از آن آتش نیز وجود ندارد؛ و اگر کسی از سطحی عبور نکرده باشد، ردپای او در آن سطح باقی نمی‌ماند. البته ممکن است ردپا یا دود ساختگی هم تولید شود، اما همین امکان نشان می‌دهد که تفسیر نمایه باید با زمینه و شواهد دیگر سنجیده شود.

\subsubsection{نماد: نشانه از راه قرارداد}

نماد یا \lr{\textit{Symbol}} نشانه‌ای است که رابطه‌ی آن با موضوع، عمدتاً از راه قرارداد، قاعده یا عادت آموخته‌شده برقرار می‌شود. واژه‌ی «درخت» شبیه درخت نیست و رابطه‌ی فیزیکیِ علّی با درخت ندارد. ما آن را به‌عنوان نام درخت می‌شناسیم، چون در زبان فارسی چنین قراردادی را آموخته‌ایم. ارقام، حروف، علائم راهنمایی و بخش بزرگی از نشانه‌های ریاضی نیز به همین شکل عمل می‌کنند.

برای مثال، علامت «\lr{+}» به‌خودی‌خود چیزی را جمع نمی‌کند. این شکل در نظام ریاضی ما قراردادی برای نشان دادن عمل جمع است. اگر این علامت را به فردی نشان دهیم که آن قرارداد را نیاموخته باشد، ممکن است صرفاً دو خط متقاطع ببیند. پس چیزی که برای ما یک نماد ریاضی است، برای او شاید فقط شکلی روی کاغذ باشد.

\begin{figure}[htbp]
\centering
\begin{tikzpicture}[node distance=2.1cm, every node/.style={draw,rounded corners,align=center,minimum width=3.0cm,minimum height=1.0cm}, >=Latex]
  \node (sign) {نشانه\\\lr{\small Sign}};
  \node[right=of sign] (obj) {موضوع\\\lr{\small Object}};
  \node[below=1.6cm of $(sign)!0.5!(obj)$] (int) {تفسیر\\\lr{\small Interpretant}};
  \draw[->,thick] (sign) -- node[above,draw=none,fill=white] {به آن مربوط است} (obj);
  \draw[->,thick] (sign) -- node[left,draw=none,fill=white] {تفسیر می‌شود} (int);
  \draw[->,thick] (obj) -- node[right,draw=none,fill=white] {در تفسیر نقش دارد} (int);
\end{tikzpicture}
\caption{در نگاه پیرس، نشانه در رابطه‌ای سه‌گانه با موضوع و تفسیر معنا پیدا می‌کند. این نمودار تنها یک نمایش ساده‌شده از رابطه است.}
\label{fig:peirce-triad}
\end{figure}

این سه دسته را نباید سه جعبه‌ی کاملاً جدا و ناسازگار تصور کنیم. یک نشانه می‌تواند در جنبه‌های مختلف، بیش از یک ویژگی داشته باشد. برای مثال، یک نقشه ممکن است از نظر شکل کلی شمایلی باشد، اما رنگ‌ها و علائم آن از راه قرارداد نمادین عمل کنند. همچنین یک عکس می‌تواند شمایلی از چهره باشد، درحالی‌که به‌عنوان اثر نوریِ ناشی از حضور همان شخص، جنبه‌ای نمایه‌ای نیز داشته باشد.

اهمیت این تقسیم‌بندی برای بحث ما در این است که «نمایندگی» فقط به شباهت وابسته نیست. گاهی یک وضعیت فیزیکی به‌خاطر شباهت، گاهی به‌خاطر ارتباط واقعی، و گاهی به‌خاطر قرارداد نقش نشانه پیدا می‌کند. بنابراین، وقتی در ادامه از سیگنال و حامل اطلاعات حرف می‌زنیم، باید همواره بپرسیم این وضعیت فیزیکی به چه چیزی مربوط شده و گیرنده با چه قاعده‌ای آن را می‌خواند.

\subsection{سیگنال چیست؟ اطلاعات در فاصله‌ی میان نشانه و فیزیک}
\label{sub:what-is-signal}

حالا می‌توانیم به واژه‌ای برسیم که در عنوان این فصل نیز حضور دارد: سیگنال. این واژه را در زندگی روزمره زیاد به کار می‌بریم. از سیگنال تلفن همراه حرف می‌زنیم، می‌گوییم دستگاه سیگنال دریافت نمی‌کند، یا از چراغ راهنمایی به‌عنوان سیگنالی برای ایستادن یاد می‌کنیم. اما از دید فیزیک و مهندسی، سیگنال دقیقاً چیست؟

به‌طور کلی، سیگنال را می‌توان کمیتی فیزیکی دانست که در زمان، مکان یا هر دو تغییر می‌کند و می‌تواند برای انتقال یا نمایش اطلاعات به کار گرفته شود. این کمیت ممکن است فشار هوا، شدت نور، ولتاژ، جریان، میدان الکترومغناطیسی یا حتی موقعیت یک قطعه‌ی مکانیکی باشد. آنچه اهمیت دارد این است که تغییرات آن کمیت بتواند در یک سامانه‌ی مشخص، از هم قابل تشخیص و تفسیر باشد. \autocite{pierce1980,shannon1948}

برای مثال، وقتی صحبت می‌کنیم، تارهای صوتی و اندام‌های گفتاری ما تغییراتی در فشار هوا ایجاد می‌کنند. این تغییرات در فضا منتشر می‌شوند و به گوش شنونده می‌رسند. گوش و مغز شنونده از این تغییرات، الگوهایی استخراج می‌کنند که به واج‌ها، واژه‌ها و در نهایت جمله‌ها مربوط می‌شوند. در اینجا، پیام همان جمله‌ای است که قصد بیانش را داریم؛ سیگنال، تغییرات فیزیکی فشار هواست؛ و کانال، محیطی است که این موج صوتی از آن عبور می‌کند.

در یک سامانه‌ی الکترونیکی نیز ممکن است ولتاژ در طول زمان تغییر کند. گیرنده با اندازه‌گیری این تغییرات و مقایسه‌ی آن‌ها با قراردادهای از پیش تعیین‌شده، می‌تواند پیام را بازسازی کند. اما ولتاژ به‌خودی‌خود جمله یا عدد نیست؛ همان‌طور که موج صوتی هم به‌خودی‌خود معنای واژه‌ها را در خود ندارد. رابطه‌ی میان سیگنال و پیام، به کدگذاری و تفسیر سامانه وابسته است.

پس بهتر است چهار مفهوم را از هم جدا نگه داریم:

\begin{itemize}\renewcommand{\labelitemi}{\lr{--}}
  \item \textbf{پیام:} محتوایی که قرار است منتقل شود؛ مثلاً عدد ۵ یا جمله‌ی «قطار رسید».
  \item \textbf{کد:} قاعده‌ای که تعیین می‌کند هر پیام یا نماد چگونه نمایش داده شود.
  \item \textbf{سیگنال:} تغییر فیزیکی‌ای که در عمل برای حمل یا نمایش پیام به کار می‌رود.
  \item \textbf{کانال:} محیط یا مسیری که سیگنال از آن عبور می‌کند.
\end{itemize}

این تمایزها در مثال‌های ساده گاهی به‌نظر بدیهی می‌رسند، اما در سامانه‌های پیچیده به‌سرعت با هم مخلوط می‌شوند. مثلاً ممکن است بگوییم «اطلاعات از طریق کابل فرستاده شد»، درحالی‌که آنچه واقعاً در کابل رخ داده، تغییرات فیزیکی ولتاژ و جریان بوده است. پیام در قالبی قراردادی به این تغییرات نسبت داده شده و گیرنده با خواندن آن تغییرات، پیام را بازیابی کرده است.

در نتیجه، سیگنال را نباید با اطلاعات یکی دانست. سیگنال چیزی است که در جهان فیزیکی رخ می‌دهد و قابل اندازه‌گیری است؛ اطلاعات کمیتی است که با توجه به یک منبع، یک احتمال، یک کد یا یک رابطه‌ی ارتباطی توصیف می‌کنیم. این دو به هم مربوط‌اند، اما یکی نیستند.

\subsection{حلقه‌های رشد درخت: ردّی از شرایط گذشته}
\label{sub:tree-rings}

برای دیدن اینکه طبیعت چگونه می‌تواند اثری از گذشته را در خود نگه دارد، لازم نیست سراغ دستگاه‌های پیچیده برویم. کافی است برش عرضی تنه‌ی یک درخت را نگاه کنیم. در چنین برشی معمولاً حلقه‌هایی را می‌بینیم که از نزدیکی مرکز تنه به سمت پوست گسترش پیدا کرده‌اند. این حلقه‌ها صرفاً طرح‌هایی تزئینی نیستند؛ در بسیاری از گونه‌های درختی که رشد فصلی مشخصی دارند، هر حلقه با یک دوره‌ی رشد سالانه ارتباط دارد.

این حلقه‌ها چگونه به وجود می‌آیند؟ درخت در طول فصل رشد، بافت چوبی تازه‌ای تولید می‌کند. با تغییر شرایط فصلی، سرعت رشد و ویژگی سلول‌هایی که ساخته می‌شوند نیز تغییر می‌کند. در بسیاری از درختان مناطق معتدل، چوبی که در ابتدای فصل رشد تشکیل می‌شود از نظر اندازه‌ی سلول‌ها و ویژگی‌های بافتی با چوبی که در ادامه‌ی فصل ساخته می‌شود تفاوت دارد. مرز میان این بخش‌ها می‌تواند به شکل حلقه‌ای قابل مشاهده باشد. درخت در فصل نامساعد ممکن است رشد بسیار کندی داشته باشد یا وارد دوره‌ی خواب شود و با آغاز فصل بعد، رشد دوباره ادامه پیدا کند. به همین دلیل درختانِ دارای فصل‌های رشد و خواب مشخص، اغلب الگویی از حلقه‌های متوالی در تنه‌ی خود باقی می‌گذارند. \autocite{frank2022dendro}

عرض این حلقه‌ها همیشه یکسان نیست. در شرایطی که عامل محدودکننده‌ی رشد در دسترس باشد، ممکن است حلقه‌ی پهن‌تری شکل بگیرد؛ در شرایط تنش، کمبود آب، دمای نامناسب یا آسیب زیستی، رشد شعاعی می‌تواند کاهش پیدا کند و حلقه باریک‌تر شود. اما همین‌جا باید مراقب یک برداشت ساده‌انگارانه باشیم: حلقه‌ی پهن الزاماً به معنای «سال پرباران» و حلقه‌ی باریک الزاماً به معنای «سال خشک» نیست. در یک منطقه ممکن است آب عامل اصلی محدودکننده‌ی رشد باشد، در منطقه‌ای دیگر دما نقش مهم‌تری داشته باشد. سن و اندازه‌ی درخت، رقابت با درختان اطراف، ترکیب خاک، نور، آفت‌ها، آتش‌سوزی و شرایط سال‌های پیشین نیز می‌توانند بر رشد اثر بگذارند. \autocite{frank2022dendro,buentgen2021treerings}

پس وقتی یک حلقه‌ی پهن یا باریک را روی چوب می‌بینیم، در واقع با یک پیام عمدی روبه‌رو نیستیم. درخت تصمیم نگرفته که برای انسان آینده گزارشی از آب‌وهوا بنویسد. حلقه نتیجه‌ی فرایندهای زیستی و محیطی است؛ اما چون شرایط محیطی بر رشد اثر گذاشته‌اند، بخشی از آن شرایط می‌تواند در ویژگی‌های چوب باقی بماند. این دقیقاً تفاوت میان «اثر فیزیکی باقی‌مانده» و «پیامی که فرستنده آگاهانه رمزگذاری کرده» است. حلقه‌ها نشانه‌هایی از گذشته‌اند، اما نه به معنای پیام‌های زبانی یا قراردادی.

برای آنکه این آثار به داده‌ی قابل استفاده تبدیل شوند، پژوهشگر باید آن‌ها را اندازه‌گیری و با شواهد دیگر مقایسه کند. درخت‌شناسان و اقلیم‌شناسان می‌توانند پهنای حلقه‌ها، چگالی چوب، ویژگی‌های سلولی یا ترکیب ایزوتوپی آن را بررسی کنند. همچنین الگوهای پهن و باریک شدن حلقه‌ها در چند درخت از یک منطقه با یکدیگر مقایسه می‌شوند تا ترتیب زمانی آن‌ها دقیق‌تر مشخص شود؛ روشی که در گاه‌شماری درختی به هم‌ترازی یا \lr{Crossdating} معروف است. سپس این اندازه‌گیری‌ها با داده‌های اقلیمیِ دوره‌ی معاصر سنجیده می‌شوند تا روشن شود در آن منطقه کدام ویژگی محیطی، چه رابطه‌ای با رشد درخت دارد. \autocite{frank2022dendro}

بنابراین حلقه‌های رشد را نباید مانند یک دماسنج یا باران‌سنج مستقیم تصور کنیم. آن‌ها آرشیوهای طبیعی‌اند که برای تفسیرشان به نمونه‌های متعدد، روش اندازه‌گیری، مقایسه‌ی منطقه‌ای و مدل‌های آماری نیاز داریم. حتی انتخاب درخت‌ها و روش تبدیل اندازه‌گیری‌ها به بازسازی اقلیم می‌تواند نتیجه را تغییر دهد. پژوهش‌های بازسازی اقلیم با حلقه‌های درختی نشان داده‌اند که تصمیم‌های تحلیلی و روش‌های آماری بر برآورد نهایی اثر می‌گذارند؛ به همین دلیل نتیجه‌ها باید همراه با عدم‌قطعیت و محدودیت‌های روش‌شناختی خوانده شوند. \autocite{buentgen2021treerings}

\begin{figure}[htbp]
\centering
\begin{tikzpicture}[scale=0.82, every node/.style={font=\small}]
  \fill[brown!12] (0,0) circle (2.45);
  \foreach \r in {2.32,2.05,1.72,1.55,1.12,0.86,0.70,0.28} {
    \draw[brown!65!black,line width=0.75pt] (0,0) circle (\r);
  }
  \fill[brown!65!black] (0,0) circle (0.055);
  \draw[->,thick] (1.83,0.68) -- (3.0,1.15) node[right,align=left,text width=2.1cm] {فاصله‌ی بیشتر میان دو حلقه: حلقه‌ی پهن‌تر};
  \draw[->,thick] (0.54,-0.54) -- (2.65,-1.35) node[right,align=left,text width=2.2cm] {فاصله‌ی کمتر میان دو حلقه: حلقه‌ی باریک‌تر};
  \node[below=2.7cm] at (0,0) {مرکز تنه};
\end{tikzpicture}
\caption{طرح شماتیک از حلقه‌های رشد در برش عرضی تنه‌ی درخت. فاصله‌ی میان مرزهای حلقه‌ها می‌تواند متفاوت باشد؛ شکل، داده‌ی اندازه‌گیری‌شده از یک درخت واقعی نیست.}
\label{fig:tree-rings-schematic}
\end{figure}

برای دیدن مسیر استنتاج، می‌توانیم آن را این‌طور دنبال کنیم: شرایط محیطی بر رشد درخت اثر می‌گذارند؛ رشد بر ویژگی‌های بافت چوب اثر می‌گذارد؛ ویژگی‌های چوب در حلقه‌ها و دیگر پارامترهای قابل اندازه‌گیری باقی می‌مانند؛ و پژوهشگر با مقایسه و کالیبراسیون آن‌ها، درباره‌ی شرایط گذشته استنتاج می‌کند. در هر گام، امکان خطا یا اثر عوامل دیگر وجود دارد و به همین دلیل نتیجه‌ی نهایی یک استنتاج علمی است، نه خواندن مستقیم یک پیام پنهان.

\begin{figure}[htbp]
\centering
\begin{tikzpicture}[node distance=0.38cm, every node/.style={draw,rounded corners,align=center,minimum width=8.9cm,minimum height=0.72cm,font=\small}, >=Latex]
  \node (env) {شرایط محیطی: آب، دما، نور، خاک و تنش‌های زیستی};
  \node[below=of env] (growth) {فرایند رشد و تشکیل سلول‌های چوبی};
  \node[below=of growth] (trace) {اثر باقی‌مانده: پهنا، چگالی و ساختار حلقه‌ها};
  \node[below=of trace] (measure) {اندازه‌گیری، مقایسه و هم‌ترازی نمونه‌ها};
  \node[below=of measure] (infer) {استنتاج محتاطانه درباره‌ی شرایط گذشته};
  \draw[->,thick] (env) -- (growth);
  \draw[->,thick] (growth) -- (trace);
  \draw[->,thick] (trace) -- (measure);
  \draw[->,thick] (measure) -- (infer);
\end{tikzpicture}
\caption{از رویدادهای محیطی تا استنتاج علمی: حلقه‌ی درخت پیام عمدی نیست، بلکه اثری فیزیکی است که از راه شناخت رابطه‌های علّی می‌توان آن را تفسیر کرد.}
\label{fig:tree-ring-inference}
\end{figure}

و همین منطق را در ردپا یا فسیل نیز می‌توان دید. ردپای یک جانور، خودِ جانور نیست؛ اما شکل و عمق آن بر اثر تماس بدن جانور با رسوب یا خاک باقی مانده است. فسیل نیز می‌تواند بخشی از ساختار یک جاندار یا اثری از فعالیت آن را حفظ کند. این آثار به ما امکان می‌دهند درباره‌ی چیزی که اکنون در برابرمان حاضر نیست استنتاج کنیم. بااین‌حال، تفسیر آن‌ها بدون شناخت جنس رسوب، شرایط دفن، فرسایش و زمینه‌ی زمین‌شناختی ممکن نیست. اثر باقی‌مانده، به‌تنهایی تمام داستان را روایت نمی‌کند.

این مثال‌ها نکته‌ای مهم را روشن می‌کنند: لازم نیست هر نشانه‌ای فرستنده‌ای داشته باشد که آگاهانه قصد ارسال پیام کرده باشد. گاهی یک رویداد طبیعی وضعیتی فیزیکی بر جای می‌گذارد و موجودی دیگر، با داشتن دانش و روش مناسب، از آن وضعیت درباره‌ی گذشته چیزی می‌آموزد. اما این هم به معنای آن نیست که هر تفاوتی در ماده خودبه‌خود اطلاعات معنادار دارد؛ باید رابطه‌ای قابل دفاع میان رویداد، اثر باقی‌مانده و استنتاجی که انجام می‌دهیم وجود داشته باشد.

\subsection{اطلاعات در طبیعت: نمونه‌ای از ژنتیک}
\label{sec:information-in-nature}

تا اینجا مثال‌های ما عمدتاً ساخته‌ی انسان بودند: نوشته، صدا، نقشه، عدد و سیگنال الکترونیکی. اما آیا بازنمایی و اطلاعات فقط در سامانه‌هایی وجود دارند که انسان آن‌ها را طراحی کرده است؟ برای پاسخ، لازم نیست از طبیعت فاصله بگیریم؛ کافی است به یکی از بنیادی‌ترین سازوکارهای حیات نگاه کنیم: \lr{DNA}.

مولکول \lr{DNA} از زنجیره‌ای از واحدهای شیمیایی به نام نوکلئوتید ساخته شده است. در ساختار رایج \lr{DNA}، چهار نوع باز نیتروژنی وجود دارد: آدنین، تیمین، گوانین و سیتوزین که معمولاً با حروف \lr{A, T, G, C} نشان داده می‌شوند. ترتیب این واحدها در امتداد مولکول می‌تواند در فرایندهای زیستی نقش داشته باشد. به همین دلیل، ما می‌توانیم توالی \lr{DNA} را به‌صورت یک رشته از نمادها بنویسیم؛ برای مثال:

\[
\text{\lr{\texttt{ATGCCATGTA...}}}
\]

البته حروف \lr{A, T, G, C} خودِ مولکول نیستند؛ آن‌ها شیوه‌ی نوشتاری ما برای اشاره به چهار نوع باز شیمیایی‌اند. در خود سلول، با مولکول‌هایی واقعی، پیوندهای شیمیایی، ساختار فضایی و سازوکارهای مولکولی روبه‌رو هستیم. اما وقتی توالی را با حروف نشان می‌دهیم، می‌توانیم درباره‌ی ترتیب آن صحبت کنیم، آن را مقایسه کنیم و الگوهایش را بررسی کنیم.

اروین شرودینگر در کتاب کوتاه اما اثرگذار خود، «حیات چیست؟» که در سال ۱۹۴۴ منتشر شد، تلاش کرد از منظر فیزیک به مسئله‌ی پایداری وراثت نگاه کند. او برای توصیف ماده‌ی وراثتی از تعبیر «جامد نامتناوب» یا عبارت انگلیسی \mbox{\lr{\textit{aperiodic solid}}} استفاده کرد و بر این نکته تأکید داشت که یک ساختار مولکولی می‌تواند در عین پایداری، آرایش پیچیده و غیرتکراری داشته باشد. این بحث در شکل‌گیری پرسش‌های بعدی درباره‌ی ماهیت مولکولی وراثت تأثیرگذار بود؛ اما باید دقت کنیم که شرودینگر در آن زمان ساختار \lr{DNA} را نمی‌شناخت و کتاب او توصیف ساختار مارپیچ دوگانه نبود. ساختار دوگانه‌ی \lr{DNA} در سال ۱۹۵۳ توسط واتسون و کریک پیشنهاد شد. \autocite{schrodinger1944,watsoncrick1953}

از این مثال می‌توان فهمید که طبیعت نیز ساختارهایی دارد که می‌توان آن‌ها را به‌عنوان حامل اطلاعات مطالعه کرد. بااین‌حال، باید از یک اشتباه رایج دوری کنیم: اینکه هر توالی نامنظم یا هر ساختار پیچیده‌ای را فوراً «اطلاعات زیاد» بنامیم. برای اینکه از اطلاعات در معنای دقیق نظریه‌ی شانون صحبت کنیم، باید یک فضای امکان‌ها و یک توزیع احتمال داشته باشیم. برای صحبت از عملکرد زیستی \lr{DNA} نیز باید سازوکارهای سلولی، تفسیر مولکولی و زمینه‌ی زیستی آن را در نظر بگیریم. صرف دیدن یک رشته‌ی غیرتکراری، به‌تنهایی نه مقدار آنتروپی آن را تعیین می‌کند و نه معنای زیستی‌اش را.

در واقع، \lr{DNA} نمونه‌ی خوبی است برای دیدن تفاوت میان سه سطح: یک توالی شیمیایی فیزیکی وجود دارد؛ ما آن توالی را با یک الفبای قراردادی بازنمایی می‌کنیم؛ و درون سلول، سازوکارهای مولکولی مشخصی وجود دارند که به بخش‌هایی از این توالی نقش زیستی می‌دهند. این سه سطح به هم وابسته‌اند، اما نباید آن‌ها را یکی بگیریم.

\section{اطلاعات و حامل آن: دو چیز که یکی نیستند}
\label{sec:information-and-carrier}

تا اینجا دیدیم که یک پیام می‌تواند به شکل‌های فیزیکی متفاوتی نمایش داده شود. عدد ۵ را می‌توان با جوهر، صدا، نور یا تغییرات ولتاژ نشان داد؛ توالی \lr{DNA} نیز می‌تواند در سطحی از توصیف، به شکل رشته‌ای از حروف \lr{A, T, G, C} نوشته شود. در همه‌ی این مثال‌ها، یک پرسش مهم باقی می‌ماند: اگر محتوای اطلاعاتی با حامل فیزیکی یکی نیست، پس رابطه‌ی دقیق میان آن‌ها چیست؟

برای اینکه مسئله روشن‌تر شود، یک برگ کاغذ را در نظر بگیرید که روی آن عدد ۵ نوشته شده است. اگر کاغذ را بسوزانیم، وضعیت فیزیکی‌ای که آن عدد را حمل می‌کرد از بین می‌رود. اما آیا خودِ عدد ۵ نیز از بین رفته است؟ واضح است که نه؛ می‌توانیم آن را دوباره روی کاغذ دیگری بنویسیم، با صدا بیان کنیم یا روی نمایشگر نشان دهیم. چیزی که از بین رفته، این نمونه‌ی خاص از ثبت فیزیکی عدد است، نه مفهوم ریاضی پنج.

اما از این مثال نباید نتیجه بگیریم که داده‌ی ثبت‌شده در جهان واقعی می‌تواند بدون هیچ حامل یا تحقق فیزیکی وجود داشته باشد. آنچه می‌توانیم بگوییم این است که یک محتوای انتزاعی ممکن است در حامل‌های متفاوتی بازنمایی شود. اگر قرار باشد آن محتوا در جایی ثبت شود، منتقل شود یا در یک سامانه‌ی فیزیکی پردازش شود، باید تحقق مشخصی داشته باشد.

\subsection{حامل، پیام و نمونه‌ی فیزیکی}
\label{sub:carrier-message-token}

در بحث‌های اطلاعاتی، گاهی لازم است سه سطح را از یکدیگر جدا کنیم. نخست، محتوای پیام است: چیزی که می‌خواهیم منتقل کنیم. دوم، کد یا قاعده‌ی بازنمایی است: روشی که تعیین می‌کند پیام چگونه نمایش داده شود. سوم، نمونه‌ی فیزیکی آن بازنمایی است: چیزی که در یک موقعیت مشخص، واقعاً وجود دارد و می‌توان آن را اندازه‌گیری یا دریافت کرد.

برای مثال، جمله‌ی «در بسته است» یک محتوای زبانی دارد. اگر آن را با خط فارسی روی کاغذ بنویسیم، کاغذ و جوهر نمونه‌ی فیزیکی آن نوشته‌اند. اگر همین جمله را با صدای بلند بگوییم، موج فشار هوا نمونه‌ی فیزیکی دیگری از همان پیام است. اگر آن را در یک فایل متنی ذخیره کنیم، الگوی مشخصی از داده در رسانه‌ی ذخیره‌سازی و وضعیت‌های داخلی دستگاه وجود دارد. در هر سه مورد، محتوای زبانی را می‌توان یکسان دانست، اما نمونه‌ی فیزیکی و شیوه‌ی بازنمایی متفاوت است.

در اینجا واژه‌ی «حامل» را برای چیزی به کار می‌بریم که امکان ثبت یا انتقال بازنمایی را در یک بستر فیزیکی فراهم می‌کند. کاغذ، موج صوتی، نور، میدان مغناطیسی و وضعیت‌های الکتریکی می‌توانند در زمینه‌های متفاوت حامل باشند. اما حامل با پیام یکی نیست. جوهر روی کاغذ به‌تنهایی نمی‌گوید چه جمله‌ای نوشته شده است؛ این کار به شکل حروف، قرارداد زبانی و توانایی خواننده برای تفسیر آن‌ها وابسته است.

همچنین نباید تصور کنیم هر چیزی که حامل یک پیام است، برای همیشه همان پیام را حمل خواهد کرد. کاغذ ممکن است خیس شود، صدا در نویز گم شود، نور پراکنده شود یا وضعیت فیزیکی یک رسانه تغییر کند. پس بازنمایی فیزیکی همیشه با مسئله‌ی پایداری، تشخیص‌پذیری و قابلیت بازیابی روبه‌روست. یک حامل خوب فقط چیزی نیست که بتواند حالت فیزیکی‌ای را ایجاد کند؛ باید بتوان آن حالت را از حالت‌های دیگر تشخیص داد و در صورت نیاز دوباره خواند.

\subsection{افلاطون‌گرایی در برابر فیزیکالیسم: آیا اطلاعات مستقل از حامل وجود دارد؟}
\label{sub:platonism-versus-physicalism}

اینجا به پرسشی می‌رسیم که از سطح مهندسی فراتر می‌رود و به فلسفه مربوط می‌شود. آیا اطلاعات چیزی مستقل از جهان فیزیکی است؟ آیا عدد پنج، رابطه‌ی ریاضی یا ساختار منطقی یک پیام، حتی اگر هیچ انسانی آن را ننویسد و هیچ دستگاهی آن را ذخیره نکند، به نوعی وجود دارد؟ یا هر آنچه واقعاً وجود دارد همان وضعیت‌های فیزیکی جهان است و اطلاعات فقط روشی است که ما برای توصیف بعضی از آن وضعیت‌ها به کار می‌بریم؟

در یک سوی این بحث می‌توان به خانواده‌ای از دیدگاه‌های افلاطون‌گرا اشاره کرد. در صورت‌بندی‌های مختلف این دیدگاه، اشیای ریاضی یا ساختارهای انتزاعی صرفاً اختراع ذهن انسان نیستند؛ آن‌ها نوعی وجود مستقل دارند و ما با استدلال و ریاضیات آن‌ها را کشف می‌کنیم. از این منظر، عدد ۵ به این دلیل وجود ندارد که روی کاغذ نوشته شده یا در حافظه‌ی رایانه ثبت شده است. نمادهای فیزیکی فقط راه‌هایی هستند که ما برای اشاره به آن ساختار انتزاعی به کار می‌بریم. البته دیدگاه‌های افلاطون‌گرا یکسان نیستند و نباید همه‌ی آن‌ها را به یک ادعای ساده فروکاست. \autocite{sep-platonism-metaphysics}

در سوی دیگر، فیزیکالیسم به‌طور کلی بر این ایده تأکید می‌کند که واقعیت از امور فیزیکی تشکیل شده یا دست‌کم هر امر واقعی به نحوی در جهان فیزیکی تحقق دارد. یک فیزیکالیست می‌تواند بگوید آنچه در جهان وجود دارد، آرایش‌های فیزیکی و روابط میان آن‌هاست؛ اطلاعات نیز یا توصیفی انتزاعی از این آرایش‌هاست یا خاصیتی است که تنها در یک سامانه‌ی فیزیکی معین معنا پیدا می‌کند. اما این هم یک خانواده‌ی گسترده از دیدگاه‌هاست و هر فیزیکالیستی الزاماً درباره‌ی جایگاه اشیای ریاضی یا ماهیت اطلاعات، پاسخ واحدی نمی‌دهد. \autocite{floridi2010,sep-physicalism,sep-multiple-realizability}

برای بحث این کتاب، لازم نیست همین‌جا یکی از این مواضع فلسفی را به‌عنوان پاسخ نهایی انتخاب کنیم. هدف ما محدودتر و در عین حال مهم‌تر است: وقتی از داده‌ای حرف می‌زنیم که در یک دستگاه ذخیره شده یا از یک کانال عبور می‌کند، درباره‌ی یک فرایند فیزیکی واقعی صحبت می‌کنیم. ممکن است محتوای ریاضی یا معناییِ آن داده را در سطحی انتزاعی توصیف کنیم، اما دستگاه برای ذخیره و بازیابی آن به حالت‌های فیزیکی قابل تمایز نیاز دارد.

این تمایز میان «وجود انتزاعی یک ساختار» و «تحقق فیزیکی یک نمونه از آن» را در مثال عدد ۵ هم دیدیم. مفهوم عدد پنج با یک لکه‌ی جوهر یکی نیست؛ اما اگر بخواهیم این عدد را روی کاغذ ثبت کنیم، باید تفاوتی فیزیکی میان حالتِ نوشته‌شده‌ی ۵ و مثلاً ۶ یا ۸ وجود داشته باشد. اگر هیچ تفاوت قابل تشخیصی میان این حالت‌ها وجود نداشته باشد، سامانه نمی‌تواند آن‌ها را از یکدیگر بازیابی کند.

\subsection{حالت منطقی و حالت فیزیکی}
\label{sub:logical-and-physical-state}

در علوم رایانه و نظریه‌ی اطلاعات، زیاد با مفهوم بیت روبه‌رو می‌شویم. بیت را معمولاً واحدی برای نمایش اطلاعات می‌دانیم که می‌تواند یکی از دو مقدار صفر یا یک را داشته باشد. اما باید دقت کنیم که صفر و یک، در این سطح، برچسب‌های انتزاعی‌اند؛ آن‌ها الزاماً به معنی نبودن و بودن الکترون، خاموش بودن و روشن بودن یک سیم، یا وجود و نبودن انرژی نیستند.

فرض کنید دستگاهی داریم که دو وضعیت قابل تشخیص دارد. می‌توانیم یکی را با برچسب ۰ و دیگری را با برچسب ۱ مشخص کنیم. اینکه این دو وضعیت دقیقاً چه هستند، به فناوری و طراحی دستگاه بستگی دارد. در یک سامانه ممکن است تفاوت میان دو سطح ولتاژ باشد؛ در سامانه‌ای دیگر ممکن است جهت مغناطش یک ناحیه، حالت نوری، یا وضعیت مکانیکی باشد. آنچه اهمیت دارد این است که دستگاه بتواند این دو وضعیت را تولید کند، نگه دارد یا تشخیص دهد.

پس باید میان حالت منطقی و حالت فیزیکی تمایز بگذاریم. حالت منطقی همان صفر یا یک است که در سطح مدل انتزاعی با آن کار می‌کنیم. حالت فیزیکی وضعیت مشخصی از یک سامانه‌ی واقعی است که آن مقدار منطقی به آن نسبت داده می‌شود. نگاشت میان این دو، بخشی از طراحی سامانه است.

\begin{figure}[htbp]
\centering
\begin{tikzpicture}[node distance=1.3cm, every node/.style={draw,rounded corners,align=center,minimum width=3.4cm,minimum height=0.95cm}, >=Latex]
  \node (logic0) {حالت منطقی ۰};
  \node[below=0.5cm of logic0] (logic1) {حالت منطقی ۱};
  \node[right=3.1cm of logic0] (phys0) {حالت فیزیکی الف};
  \node[below=0.5cm of phys0] (phys1) {حالت فیزیکی ب};
  \draw[->,thick] (logic0) -- node[above,draw=none,fill=white] {نگاشت} (phys0);
  \draw[->,thick] (logic1) -- node[below,draw=none,fill=white] {نگاشت} (phys1);
  \node[draw=none,below=0.5cm of $(logic1)!0.5!(phys1)$,text width=8cm] {یک مقدار منطقی می‌تواند با فناوری‌های فیزیکی متفاوتی پیاده‌سازی شود.};
\end{tikzpicture}
\caption{مقادیر منطقی صفر و یک، برچسب‌هایی انتزاعی‌اند که به حالت‌های فیزیکی قابل تشخیص نگاشت می‌شوند.}
\label{fig:logical-physical-state}
\end{figure}

این موضوع را می‌توانیم از زاویه‌ای دیگر هم ببینیم. یک ساعت عقربه‌ای زمان را با موقعیت زاویه‌ای عقربه نمایش می‌دهد؛ یک ساعت دیجیتال آن را با رقم‌هایی روی صفحه نشان می‌دهد. هر دو می‌توانند زمان را نمایش دهند، اما ساختار فیزیکی و روش خواندنشان متفاوت است. پس لازم نیست یک محتوای منطقی به یک ماده یا سازوکار فیزیکی خاص وابسته باشد. کافی است سامانه بتواند میان حالت‌های مرتبط تمایز بگذارد و آن‌ها را به‌درستی بخواند.

همین نکته بعدها در بحث پیاده‌سازی چندگانه نیز اهمیت پیدا می‌کند: یک فرایند یا محتوای انتزاعی ممکن است در بسترهای فیزیکی مختلف تحقق پیدا کند. این به معنای بی‌نیازی آن از فیزیک نیست؛ بلکه به این معناست که یک نوع تحقق فیزیکیِ یگانه و اجباری برای آن وجود ندارد. \autocite{putnam1967,fodor1974,sep-multiple-realizability}

\subsection{شیطان ماکسول: آیا اطلاعات می‌تواند بر جهان فیزیکی اثر بگذارد؟}
\label{sub:maxwells-demon}

تا اینجا بیشتر از این صحبت کردیم که جهان فیزیکی چگونه می‌تواند حامل اطلاعات باشد. اما آیا رابطه فقط یک‌طرفه است؟ آیا اطلاعات و دانستن چیزی درباره‌ی یک سامانه می‌تواند در کار فیزیکی آن سامانه نقشی داشته باشد؟ یکی از مشهورترین مثال‌های تاریخ فیزیک برای طرح این پرسش، آزمایش ذهنی «شیطان ماکسول» است.

جیمز کلرک ماکسول در قرن نوزدهم آزمایش ذهنی‌ای را مطرح کرد که بعدها به یکی از بحث‌های مهم درباره‌ی رابطه‌ی اطلاعات و ترمودینامیک تبدیل شد. یک جعبه‌ی بسته را تصور کنید که با دیواره‌ای به دو بخش تقسیم شده است. در این دیواره دریچه‌ی کوچکی وجود دارد و موجودی بسیار کوچک و فرضی می‌تواند با باز و بسته کردن آن، مولکول‌های گاز را از یک سمت به سمت دیگر عبور دهد. این موجود را «شیطان» نامیدند؛ نه به معنای موجودی فراطبیعی، بلکه فقط برای اینکه نقش ناظری فرضی را بازی کند که می‌تواند سرعت مولکول‌ها را تشخیص دهد. این آزمایش ذهنی در بحث ماکسول درباره‌ی نظریه‌ی جنبشی گرما مطرح شد. \autocite{maxwell1871}

در یک گاز معمولی، مولکول‌ها با سرعت‌های گوناگون حرکت می‌کنند و در حالت تعادل، تفاوت دمای پایدار میان دو بخش جعبه باقی نمی‌ماند. اما شیطان می‌تواند مولکول‌های سریع‌تر را به یک سمت و مولکول‌های کندتر را به سمت دیگر هدایت کند. در نتیجه، به نظر می‌رسد می‌توان بدون صرف کار، دو بخش با دماهای متفاوت ساخت؛ گویی کاهش آنتروپی رخ داده و قانون دوم ترمودینامیک نقض شده است.

اما شیطان چگونه تصمیم می‌گیرد کدام مولکول سریع‌تر است؟ باید وضعیت مولکول‌ها را اندازه‌گیری کند، اطلاعاتی درباره‌ی آن‌ها به دست آورد و این اطلاعات را برای تصمیم‌گیری نگه دارد. به این ترتیب، مسئله‌ی فیزیکی به مسئله‌ی اطلاعاتی گره می‌خورد. شیطان صرفاً دریچه را باز و بسته نمی‌کند؛ باید چیزی را بداند و بر اساس آن دانسته عمل کند.

در قرن بیستم، لئو زیلارد در مدلی ساده‌تر نشان داد که رابطه‌ی میان اندازه‌گیری، اطلاعات و ترمودینامیک را می‌توان به شکلی دقیق‌تر بررسی کرد. در موتور تک‌مولکولی مشهور او، دانستن اینکه مولکول در کدام سمت محفظه قرار دارد می‌تواند امکان استخراج کار از یک چرخه‌ی کنترل‌شده را فراهم کند. اما چرخه‌ی کامل، از جمله آماده‌سازی و بازنشانی حافظه‌ی عامل اندازه‌گیرنده، باید در تحلیل وارد شود. نمی‌توان فقط مرحله‌ای را ببینیم که در آن اطلاعات به دست آمده و از آن کار استخراج شده است و هزینه‌ی باقی فرایند را نادیده بگیریم. \autocite{szilard1929,leffrex2003}

\subsection{اصل لاندائر: پاک کردن اطلاعات چه هزینه‌ای دارد؟}
\label{sub:landauer-principle}

رولف لاندائر در سال ۱۹۶۱ پیوند مهمی میان محاسبه و ترمودینامیک مطرح کرد. ایده‌ی اصلی او این بود که محاسبه را نباید چیزی کاملاً جدا از جهان فیزیکی بدانیم. هر محاسبه‌ای روی یک سامانه‌ی واقعی اجرا می‌شود و سامانه‌ی واقعی نیز تابع محدودیت‌های فیزیکی است. اما از این حرف نباید نتیجه بگیریم که هر بار یک بیت جابه‌جا می‌شود، الزاماً مقدار ثابتی انرژی مصرف یا گرما تولید می‌شود. نتیجه‌ی دقیق‌تر لاندائر درباره‌ی عملیات منطقیِ برگشت‌ناپذیر، به‌ویژه پاک کردن اطلاعات، است. \autocite{landauer1961}

برای فهم این نکته، یک حافظه‌ی فرضی را تصور کنید که در ابتدا ممکن است صفر یا یک را نگه داشته باشد. اگر بخواهیم آن را به‌گونه‌ای بازنشانی کنیم که در پایان حتماً صفر باشد، باید هر دو حالت آغازین را به یک حالت نهایی واحد تبدیل کنیم:

\[
0\longrightarrow 0,\qquad 1\longrightarrow 0
\]

در این فرایند، تفاوت میان دو حالت آغازین از بین می‌رود. پیش از پاک‌سازی، سامانه می‌توانسته دو حالت ممکن را از هم تشخیص دهد؛ پس از آن، دیگر نمی‌توان فهمید که ابتدا کدام حالت وجود داشته است. این نگاشت چندبه‌یک، برگشت‌ناپذیر است، چون از روی خروجی نهایی نمی‌توان ورودی اولیه را بازسازی کرد.

اصل لاندائر بیان می‌کند که در شرایط استاندارد، پاک کردن یک بیتِ نامعلوم در دمای مطلق \(T\)، دست‌کم با پراکندگی گرمایی برابر با مقدار زیر همراه است:

\[
Q_{\min}=k_{\mathrm B}T\ln 2
\]

در این رابطه، \(k_{\mathrm B}\) ثابت بولتزمن و \(T\) دمای محیط است. این مقدار یک حد پایین نظری است، نه پیش‌بینی مقدار گرمایی که هر رایانه‌ی امروزی در هر عملیات واقعاً تولید می‌کند. دستگاه‌های واقعی معمولاً بسیار بیشتر از این حد انرژی مصرف می‌کنند، چون اجزای فیزیکی، مقاومت‌ها، خازن‌ها، نشتی‌ها، تغییرات دما و محدودیت‌های ساخت در کارند. \autocite{landauer1961}

نکته‌ی اساسی این است که اصل لاندائر درباره‌ی هر نوع پردازش اطلاعات نیست. عملیات منطقی برگشت‌پذیر، دست‌کم در اصل و در شرایط ایده‌آل، می‌تواند بدون این حداقل هزینه‌ی اجتناب‌ناپذیرِ پاک‌سازی انجام شود. آنچه حد لاندائر را فعال می‌کند، از دست رفتن برگشت‌ناپذیر اطلاعات منطقی است؛ نمونه‌ی ساده‌ی آن پاک کردن یک بیت نامعلوم است. بنابراین جمله‌ی «هر پردازش اطلاعاتی حداقل به اندازه‌ی \(k_{\mathrm B}T\ln2\) انرژی مصرف می‌کند» نادرست است. صورت دقیق‌تر این است که پاک کردن برگشت‌ناپذیر یک بیت اطلاعات، تحت فرض‌های استاندارد، حداقل این مقدار گرما را به محیط منتقل می‌کند. \autocite{bennett1982,georgescu2021landauer}

\begin{figure}[htbp]
\centering
\begin{tikzpicture}[node distance=1.2cm, every node/.style={draw,rounded corners,align=center,minimum width=2.4cm,minimum height=0.9cm}, >=Latex]
  \node (in0) {\lr{0}};
  \node[below=0.55cm of in0] (in1) {\lr{1}};
  \node[right=3.8cm of in0,yshift=-0.55cm] (out) {\lr{0}};
  \draw[->,thick] (in0.east) -- (out.west);
  \draw[->,thick] (in1.east) -- (out.west);
  \node[draw=none,below=0.75cm of out,text width=4.8cm] {دو حالت آغازین به یک حالت نهایی نگاشت می‌شوند؛ تمایز اولیه از دست می‌رود.};
\end{tikzpicture}
\caption{نمونه‌ای ساده از پاک‌سازی یک بیت: دو حالت ممکن به یک حالت نهایی تبدیل می‌شوند.}
\label{fig:landauer-erasure}
\end{figure}

حالا می‌توانیم به شیطان ماکسول برگردیم. اگر موجود فرضی بخواهد دائماً اطلاعات مولکول‌ها را ثبت کند، از آن برای کنترل دریچه بهره ببرد و در پایان چرخه حافظه‌اش را دوباره به وضعیت آغازین برگرداند، بازنشانی حافظه را نمی‌توان از تحلیل حذف کرد. در روایت استاندارد لاندائر و بنت، هزینه‌ی ترمودینامیکی پاک کردن اطلاعات مانع آن می‌شود که شیطان بدون هزینه و برای همیشه قانون دوم را دور بزند. البته جزئیات فیزیکی اندازه‌گیری، حافظه، کنترل و بازنشانی همچنان موضوع پژوهش است؛ نکته‌ی مهم برای ما این است که اطلاعات در یک سامانه‌ی فیزیکی می‌تواند به منابع و محدودیت‌های فیزیکی گره بخورد. \autocite{bennett1982,georgescu2021landauer}

بنابراین، تا اینجا می‌توانیم یک نتیجه‌ی محدود اما مهم بگیریم: داده در سطح انتزاعی با حامل خود یکی نیست، اما وقتی داده در یک سامانه‌ی فیزیکی ثبت، نگهداری یا پاک می‌شود، نحوه‌ی بازنمایی آن و عملیات انجام‌شده روی آن می‌تواند پیامدهای فیزیکی داشته باشد. این هنوز به معنای آن نیست که «اطلاعات ماده است» یا هر بیت یک ذره‌ی مستقل در طبیعت است. معنای دقیق‌ترش این است که پردازش اطلاعات در جهان واقعی باید به‌وسیله‌ی فرایندهای فیزیکی انجام شود و از قوانین فیزیک بیرون نیست.

\section{از پیوسته به گسسته: چرا سیگنال دیجیتال؟}
\label{sec:continuous-to-discrete}

تا اینجا فهمیدیم که اطلاعات برای انتقال و ذخیره شدن باید در یک وضعیت فیزیکی قابل تشخیص بازنمایی شود. اما حالا یک مشکل عملی داریم: بسیاری از کمیت‌های فیزیکی در جهان به‌صورت پیوسته تغییر می‌کنند. صدای انسان، شدت نور، دمای اتاق و فشار هوا، در حالت معمول به چند مقدار محدود و جداگانه تقسیم نشده‌اند. اگر بخواهیم این پدیده‌ها را در یک سامانه‌ی دیجیتال پردازش کنیم، باید آن‌ها را به شکلی گسسته و قابل نمایش تبدیل کنیم.

برای روشن‌تر شدن مسئله، صدای یک نفر را تصور کنید. وقتی او حرف می‌زند، فشار هوا در هر لحظه تغییر می‌کند. این تغییرات را می‌توان به‌صورت یک موج ثبت کرد. موج صوتی در زمان پیوسته تغییر می‌کند و دامنه‌ی آن نیز می‌تواند در یک بازه، مقدارهای بسیار زیادی بگیرد. اما رایانه برای ذخیره و پردازش داده، معمولاً با مجموعه‌ای محدود از حالت‌های نمادین کار می‌کند. پس باید از یک نمایش پیوسته به یک نمایش گسسته برسیم.

این گذار فقط یک ترفند رایانه‌ای نیست. در هر ابزار اندازه‌گیری دیجیتال، دوربین، ضبط‌کننده‌ی صدا، اسیلوسکوپ دیجیتال یا حسگر هوشمند، با نسخه‌ای از همین مسئله روبه‌رو می‌شویم: چگونه یک پدیده‌ی پیوسته را به مجموعه‌ای از مقدارهای گسسته تبدیل کنیم، بدون آنکه اطلاعات مهم آن را بیش از حد از دست بدهیم؟

\subsection{سیگنال آنالوگ چیست؟}
\label{sub:analog-signal}

به سیگنالی آنالوگ می‌گوییم که مقدار آن می‌تواند در یک بازه به‌صورت پیوسته تغییر کند و این تغییرات، به‌نحوی متناظر با کمیتی دیگر باشند. برای مثال، میکروفون می‌تواند تغییرات فشار هوا را به تغییرات ولتاژ تبدیل کند. در خروجی میکروفون، ولتاژ در طول زمان بالا و پایین می‌رود و شکل موج آن تا حدی با تغییرات صدای ورودی متناظر است.

اگر صدای بلندتری وارد میکروفون شود، معمولاً دامنه‌ی سیگنال خروجی نیز بیشتر می‌شود؛ و اگر فرکانس صدای ورودی تغییر کند، الگوی زمانی سیگنال هم تغییر خواهد کرد. البته این رابطه دقیقاً یکسان و بی‌نقص نیست؛ میکروفون پاسخ فرکانسی، حساسیت و محدودیت‌های فیزیکی خودش را دارد. اما در اصل، تغییرات یک کمیت پیوسته به تغییرات کمیت دیگری تبدیل شده‌اند.

واژه‌ی «آنالوگ» به همین شباهت یا تناظر اشاره دارد: تغییرات سیگنال، به شکلی پیوسته تغییرات کمیت موردنظر را دنبال می‌کند. برای مثال، در یک دماسنج جیوه‌ای، ارتفاع ستون جیوه با دما تغییر می‌کند. در یک سامانه‌ی صوتی آنالوگ، ولتاژ می‌تواند تغییرات فشار هوا را بازنمایی کند. در هر دو مورد، مقدار فیزیکی خروجی در یک بازه‌ی پیوسته تغییر می‌کند.

اما این پیوستگی، به معنی دقت نامحدود نیست. ابزار اندازه‌گیری واقعی با نویز، محدودیت حساسیت، اصطکاک، پهنای باند و خطای ساخت روبه‌رو است. حتی اگر مدل ریاضی سیگنال پیوسته باشد، ابزار واقعی نمی‌تواند هر تغییر بی‌نهایت کوچکی را تشخیص دهد. بنابراین نباید «پیوسته بودن در مدل» را با «دقت بی‌نهایت در دستگاه واقعی» اشتباه بگیریم.

\subsection{نمونه‌برداری: ثبت لحظه‌های جداگانه}
\label{sub:sampling}

اولین گام برای تبدیل یک سیگنال پیوسته به داده‌ی دیجیتال، نمونه‌برداری است. به‌جای اینکه مقدار سیگنال را در تمام لحظه‌ها ثبت کنیم، در زمان‌های مشخص از آن نمونه می‌گیریم. مثلاً اگر در هر ثانیه ۴۴٬۱۰۰ بار مقدار یک سیگنال صوتی را اندازه‌گیری کنیم، نرخ نمونه‌برداری ما ۴۴٫۱ کیلوهرتز است.

فرض کنید یک موج صوتی را روی کاغذ رسم کرده‌ایم. اگر بخواهیم آن را نمونه‌برداری کنیم، می‌توانیم در فاصله‌های زمانی منظم، مقدار دامنه را بخوانیم و ثبت کنیم. نتیجه دیگر یک منحنی پیوسته نیست، بلکه مجموعه‌ای از عددهاست که هر کدام مقدار سیگنال را در یک زمان مشخص نشان می‌دهند.

اگر فاصله‌ی زمانی میان نمونه‌ها را با \(T_s\) نشان دهیم، نرخ نمونه‌برداری برابر است با:

\[
f_s=\frac{1}{T_s}
\]

یعنی هرچه فاصله‌ی زمانی میان نمونه‌ها کمتر باشد، در هر ثانیه نمونه‌های بیشتری ثبت می‌کنیم. اما آیا با افزایش تعداد نمونه‌ها می‌توانیم هر سیگنالی را به‌درستی بازسازی کنیم؟ پاسخ به ویژگی‌های فرکانسی سیگنال وابسته است.

\subsection{قضیه‌ی نمونه‌برداری نایکوئیست و شانون}
\label{sub:nyquist-shannon-sampling}

یک سیگنال صوتی یا الکترونیکی ممکن است مؤلفه‌های فرکانسی مختلفی داشته باشد. اگر سریع‌ترین تغییرات آن را نادیده بگیریم و با نرخ نمونه‌برداری بیش از حد پایین از آن نمونه بگیریم، ممکن است اطلاعات مهمی را از دست بدهیم. حتی گاهی سیگنالی که در واقع فرکانس بالایی دارد، پس از نمونه‌برداری با فرکانس پایین‌تر، به‌اشتباه به شکل سیگنالی با فرکانس دیگر دیده می‌شود. به این پدیده هم‌پوشانی طیفی یا \lr{\textit{Aliasing}} می‌گوییم.

قضیه‌ی نمونه‌برداری نایکوئیست و شانون برای سیگنال‌های محدودباند می‌گوید که اگر بالاترین فرکانس موجود در سیگنال \(f_{\max}\) باشد، برای بازسازی ایده‌آل آن باید نرخ نمونه‌برداری از دو برابر این مقدار بیشتر باشد:

\[
f_s>2f_{\max}
\]

این مقدار را نرخ نایکوئیست می‌نامند. برای مثال، اگر سیگنالی هیچ مؤلفه‌ای بالاتر از ۱۰ کیلوهرتز نداشته باشد، نرخ نمونه‌برداری باید بیش از ۲۰ هزار نمونه در ثانیه باشد تا تحت فرض‌های قضیه امکان بازسازی ایده‌آل فراهم شود. \autocite{oppenheim2010,cover2006}

اما در عمل، نمی‌توانیم فقط به این رابطه نگاه کنیم و فرض کنیم هر دستگاهی دقیقاً با دو برابر بالاترین فرکانس، سیگنال را بی‌نقص ثبت خواهد کرد. سیگنال‌های واقعی ممکن است مؤلفه‌های فرکانسی ناخواسته داشته باشند، فیلترها ایده‌آل نیستند، و نرخ نمونه‌برداری با محدودیت سخت‌افزاری روبه‌روست. به همین دلیل، پیش از نمونه‌برداری معمولاً از فیلتر پایین‌گذر استفاده می‌شود تا مؤلفه‌های فرکانسی بالاتر از محدوده‌ی موردنظر کاهش پیدا کنند. این فیلتر ضدهم‌پوشانی است و در منابع انگلیسی آن را \lr{\textit{Anti-aliasing Filter}} می‌نامند.

یک نکته‌ی تاریخی هم ارزش یادآوری دارد. هری نایکوئیست در مقاله‌ی سال ۱۹۲۴ خود درباره‌ی سرعت تلگراف، عوامل مؤثر بر سرعت ارسال پیام و شکل‌دهی سیگنال را بررسی کرد. قضیه‌ای که امروز در قالب نمونه‌برداری نایکوئیست و شانون می‌شناسیم، صورت‌بندی دقیق‌تری از شرط بازسازی سیگنال محدودباند است که در کارهای بعدی توسعه یافت. بنابراین نباید تمام محتوای قضیه‌ی نمونه‌برداری امروزی را مستقیماً به مقاله‌ی ۱۹۲۴ نایکوئیست نسبت دهیم. \autocite{nyquist1924,cover2006}

\begin{figure}[htbp]
\centering
\begin{tikzpicture}[x=0.75cm,y=1.2cm]
  \draw[->] (0,0) -- (10.8,0) node[right] {زمان};
  \draw[->] (0,-1.45) -- (0,1.7) node[above] {دامنه};
  \draw[domain=0:10,samples=150,smooth,thick,genixblue] plot (\x,{sin(\x r)});
  \foreach \x in {0,0.8,...,10} {
    \fill[genixteal] (\x,{sin(\x r)}) circle (2pt);
    \draw[genixteal,densely dotted] (\x,0) -- (\x,{sin(\x r)});
  }
\end{tikzpicture}
\caption{نمونه‌برداری از یک موج پیوسته؛ نقاط مشخص‌شده مقدار سیگنال را در لحظه‌های جداگانه نشان می‌دهند.}
\label{fig:sampling}
\end{figure}

\subsection{کوانتیزه‌سازی: از مقدار پیوسته به سطح‌های محدود}
\label{sub:quantization}

نمونه‌برداری به‌تنهایی برای ساخت یک داده‌ی دیجیتال کافی نیست. پس از آن، هنوز با مجموعه‌ای از مقدارهای اندازه‌گیری‌شده روبه‌رو هستیم که ممکن است هر کدام در یک بازه‌ی پیوسته قرار بگیرند. برای اینکه این مقدارها به شکل دیجیتال ذخیره شوند، باید آن‌ها را به مجموعه‌ای محدود از سطح‌های مجاز نگاشت کنیم. این فرایند را کوانتیزه‌سازی می‌نامیم.

فرض کنید می‌خواهیم ولتاژی را که بین صفر تا یک ولت تغییر می‌کند، با چهار بیت نمایش دهیم. چهار بیت می‌توانند ۱۶ حالت متفاوت بسازند:

\[
2^4=16
\]

پس می‌توانیم بازه‌ی صفر تا یک ولت را به ۱۶ سطح تقسیم کنیم و هر سطح را با یکی از ۱۶ کد چهار بیتی نمایش دهیم. اگر به‌جای چهار بیت از هشت بیت استفاده کنیم، ۲۵۶ سطح خواهیم داشت؛ و با شانزده بیت، ۶۵٬۵۳۶ سطح. هرچه تعداد بیت‌ها بیشتر باشد، سطح‌های قابل نمایش متراکم‌تر می‌شوند و خطای کوانتیزه‌سازی می‌تواند کاهش یابد.

البته هر اندازه‌گیری پس از کوانتیزه‌سازی ممکن است اندکی با مقدار اصلی تفاوت داشته باشد. اگر مقدار واقعی بین دو سطح مجاز قرار بگیرد، سامانه معمولاً نزدیک‌ترین سطح را انتخاب می‌کند. اختلاف میان مقدار واقعی و مقدار نمایش‌داده‌شده را خطای کوانتیزه‌سازی می‌نامیم.

این خطا را می‌توان در یک مثال ساده دید. فرض کنید مقدار واقعی یک سیگنال ۰٫۳۷ ولت باشد، اما سطح‌های مجاز دستگاه فقط ۰٫۳۰ و ۰٫۴۰ ولت را در نزدیکی آن داشته باشند. دستگاه ممکن است مقدار را به ۰٫۴۰ ولت گرد کند. در این صورت، مقدار ذخیره‌شده دقیقاً همان مقدار اولیه نیست؛ بلکه تقریب دیجیتالی آن است.

پس تبدیل آنالوگ به دیجیتال دو مرحله‌ی اصلی دارد: نمونه‌برداری در زمان، و کوانتیزه‌سازی در دامنه. اولی زمان‌های اندازه‌گیری را گسسته می‌کند؛ دومی مجموعه‌ی مقدارهای قابل نمایش را محدود می‌کند. این دو عملیات با یکدیگر متفاوت‌اند و نباید آن‌ها را یک فرایند واحد تصور کنیم.

\begin{figure}[htbp]
\centering
\begin{tikzpicture}[x=0.85cm,y=1.0cm]
  \draw[->] (0,0) -- (9.8,0) node[right] {زمان};
  \draw[->] (0,0) -- (0,3.7) node[above] {دامنه};
  \draw[step=1,gray!35,dotted] (0,0) grid (9,3.3);
  \draw[thick,genixblue] plot[smooth] coordinates {(0,0.5)(1,1.4)(2,2.3)(3,3.0)(4,2.6)(5,1.5)(6,0.8)(7,1.1)(8,2.0)(9,2.7)};
  \foreach \x/\y in {0/0.5,1/1.4,2/2.3,3/3.0,4/2.6,5/1.5,6/0.8,7/1.1,8/2.0,9/2.7} {
    \pgfmathsetmacro{\q}{round(\y)}
    \fill[genixteal] (\x,\q) circle (2.4pt);
    \draw[genixteal,dashed] (\x,\y) -- (\x,\q);
  }
\end{tikzpicture}
\caption{در کوانتیزه‌سازی، هر نمونه به نزدیک‌ترین سطح مجاز نگاشت می‌شود؛ این کار می‌تواند خطایی محدود ایجاد کند.}
\label{fig:quantization}
\end{figure}

\subsection{سیگنال دیجیتال: گسسته بودن به چه معناست؟}
\label{sub:digital-signal}

وقتی از سیگنال دیجیتال حرف می‌زنیم، منظورمان سیگنالی است که اطلاعات آن با مجموعه‌ای از نمادهای گسسته نمایش داده می‌شود. رایج‌ترین حالت در رایانه‌های امروزی استفاده از دو نماد صفر و یک است. اما این نکته را نباید فراموش کنیم که حتی در یک دستگاه دیجیتال، سیگنال در نهایت یک پدیده‌ی فیزیکی است. ولتاژ روی یک پایه‌ی تراشه، جریان در یک مسیر، نور در یک فیبر یا میدان مغناطیسی در یک رسانه، همگی کمیت‌هایی فیزیکی‌اند که در طول زمان و مکان تغییر می‌کنند.

پس دیجیتال بودن به این معنا نیست که سیگنال از فیزیک جدا شده یا دیگر مقدار فیزیکی ندارد. تفاوت در این است که اطلاعات را به سطح‌های گسسته نسبت می‌دهیم. مثلاً یک بازه‌ی ولتاژ را به‌عنوان صفر منطقی و بازه‌ای دیگر را به‌عنوان یک منطقی تفسیر می‌کنیم. دستگاه قرار نیست هر مقدار کوچک ولتاژ را یک نماد مستقل بداند؛ بلکه بازه‌های مشخصی از مقدارها را به یک حالت منطقی نسبت می‌دهد.

این روش یک مزیت مهم دارد: تا وقتی مقدار سیگنال در محدوده‌ی مجاز باقی بماند، تغییرهای کوچک و ناخواسته‌ی ولتاژ لزوماً معنای دیجیتال آن را عوض نمی‌کنند. اگر برای نمایش یک، بازه‌ای از ولتاژها مجاز باشد، اندکی نویز هنوز می‌تواند در همان بازه باقی بماند و گیرنده همچنان یک را تشخیص دهد. این حاشیه‌ی تحمل در برابر نویز را حاشیه‌ی نویز می‌نامیم.

البته این مصونیت مطلق نیست. اگر نویز یا افت سیگنال به‌اندازه‌ای زیاد باشد که مقدار از محدوده‌ی معتبر خارج شود، ممکن است یک به صفر یا برعکس خوانده شود؛ یا مقدار در محدوده‌ای قرار بگیرد که گیرنده نتواند آن را با اطمینان تفسیر کند. بنابراین دیجیتال بودن خطا را ناممکن نمی‌کند؛ بلکه با تعریف سطح‌های گسسته و قواعد تصمیم‌گیری، می‌تواند تشخیص و کنترل خطا را آسان‌تر کند.

از همین‌جا تفاوت آنالوگ و دیجیتال روشن‌تر می‌شود. در یک سامانه‌ی آنالوگ، هر تغییر کوچک در دامنه‌ی سیگنال ممکن است به تغییر متناظر در مقدار بازنمایی‌شده منجر شود. در یک سامانه‌ی دیجیتال، مقدارهای فیزیکی در بازه‌هایی دسته‌بندی می‌شوند و هر بازه به یک نماد گسسته نسبت داده می‌شود. این دسته‌بندی می‌تواند مقاومت سامانه را در برابر تغییرهای کوچک افزایش دهد، اما در عوض به انتخاب سطح‌ها، آستانه‌ها و قواعد کدگذاری نیاز دارد.

\subsection{آیا صفر و یک واقعاً دو ولتاژ مشخص‌اند؟}
\label{sub:logic-levels}

در مثال‌های مقدماتی گاهی می‌گویند صفر منطقی برابر صفر ولت و یک منطقی برابر پنج ولت است. این مثال می‌تواند برای آشنایی اولیه مفید باشد، اما نباید آن را یک قانون همگانی تصور کنیم. خانواده‌های مختلف مدارهای دیجیتال، ولتاژهای تغذیه و استانداردهای متفاوتی دارند. حتی در یک استاندارد مشخص هم معمولاً صفر و یک با یک مقدار دقیق تعریف نمی‌شوند؛ بلکه با محدوده‌ها و آستانه‌ها مشخص می‌شوند.

برای نمونه، ممکن است یک گیرنده هر ولتاژی پایین‌تر از یک آستانه را صفر منطقی و هر ولتاژی بالاتر از آستانه‌ی دیگری را یک منطقی در نظر بگیرد. میان این دو آستانه ناحیه‌ای وجود دارد که نباید آن را وضعیت معتبر و تضمین‌شده‌ی صفر یا یک تلقی کرد. این ناحیه برای جلوگیری از تصمیم‌گیری نامطمئن اهمیت دارد. مقادیر دقیق این آستانه‌ها وابسته به استاندارد، فناوری و مشخصات قطعه‌اند؛ بنابراین یک عدد واحد برای تمام دستگاه‌ها وجود ندارد. \autocite{cambridge-digital-electronics,ni-logic-levels}

در نتیجه، بیت به‌عنوان مقدار منطقی با ولتاژ یکی نیست. ولتاژ فقط یکی از راه‌های رایج پیاده‌سازی فیزیکی بیت در مدارهای الکترونیکی است. این نکته ما را به پرسشی می‌رساند که فصل را به سمت مبانی برق هدایت می‌کند: اگر صفر و یک در یک رایانه با وضعیت‌های الکتریکی بازنمایی می‌شوند، این وضعیت‌های الکتریکی دقیقاً چه هستند و چگونه می‌توانیم آن‌ها را تغییر دهیم و اندازه‌گیری کنیم؟

\section{چرا الکتریسیته؟ از بیت انتزاعی تا ولتاژ}
\label{sec:why-electricity}

اکنون به پرسشی رسیده‌ایم که از ابتدا در پسِ بحث ما قرار داشت. اگر یک بیت فقط یک مقدار انتزاعی میان صفر و یک است، چگونه در یک رایانه‌ی واقعی حضور پیدا می‌کند؟ چگونه دستگاه می‌تواند آن را ذخیره کند، تغییر دهد، بخواند یا به بخش دیگری بفرستد؟ در بخش‌های قبل دیدیم که یک مقدار منطقی باید به یک وضعیت فیزیکی نگاشت شود. حالا می‌خواهیم یکی از رایج‌ترین راه‌های این نگاشت را بررسی کنیم: استفاده از پدیده‌های الکتریکی.

شاید بپرسید چرا الکتریسیته؟ چرا رایانه‌ها از نور، حرکت مکانیکی، صدا یا هر پدیده‌ی فیزیکی دیگری استفاده نمی‌کنند؟ واقعیت این است که از این پدیده‌ها نیز استفاده می‌شود. فیبر نوری داده را با نور منتقل می‌کند؛ دیسک‌های مغناطیسی از وضعیت مغناطش بهره می‌برند؛ و سامانه‌های مکانیکی یا شیمیایی هم می‌توانند حالت‌های اطلاعاتی داشته باشند. بنابراین الکتریسیته تنها راه ممکن نیست.

اما الکتریسیته چند ویژگی دارد که آن را برای ساخت سامانه‌های محاسباتی بسیار مناسب کرده است: می‌توان اختلاف پتانسیل و جریان را در مسیرهای کوچک کنترل کرد، تغییرات آن‌ها را با سرعت بالا اندازه گرفت، آن‌ها را به‌سادگی به سیگنال‌های دیگر تبدیل کرد و در مدارهای پیچیده به کار گرفت. همچنین می‌توان بسیاری از اجزای الکترونیکی را در ابعاد بسیار کوچک کنار هم ساخت. به همین دلیل، از رایانه‌های شخصی تا سرورها و تلفن‌های همراه، بخش بزرگی از محاسبات امروزی با سامانه‌های الکترونیکی انجام می‌شود.

اما پیش از آنکه درباره‌ی یک بیت الکتریکی صحبت کنیم، باید چند مفهوم ساده را دقیق بشناسیم. ولتاژ چیست؟ جریان چیست؟ بار الکتریکی چه نقشی دارد؟ و چرا یک اختلاف ولتاژ می‌تواند به‌عنوان حامل یک نماد منطقی به کار رود؟

\subsection{بار الکتریکی و میدان الکتریکی}
\label{sub:charge-and-electric-field}

برای شروع، باید به مفهوم بار الکتریکی برسیم. ماده از ذراتی مانند پروتون‌ها و الکترون‌ها تشکیل شده است. پروتون بار الکتریکی مثبت و الکترون بار الکتریکی منفی دارد. این بارها ویژگی‌ای بنیادی از ذرات‌اند و باعث می‌شوند میان آن‌ها نیروهای الکتریکی برقرار شود. بارهای هم‌نام یکدیگر را دفع می‌کنند و بارهای ناهم‌نام یکدیگر را جذب می‌کنند.

مقدار بار الکتریکی را با نماد \(Q\) یا \(q\) نشان می‌دهیم و واحد آن در دستگاه \lr{SI} کولن است. بار یک الکترون تقریباً برابر است با:

\[
e=1.602\times10^{-19}\;\mathrm{C}
\]

علامت بار الکترون منفی است؛ در رابطه‌ی بالا اندازه‌ی بار بنیادی را نوشته‌ایم. این مقدار بسیار کوچک است، به همین دلیل حتی یک کولن بار الکتریکی معادل تعداد بسیار زیادی بار بنیادی است.

بار الکتریکی فقط در یک شیء ساکن یا یک ذره‌ی جداگانه اهمیت ندارد. بارها در اطراف خود میدان الکتریکی ایجاد می‌کنند. میدان الکتریکی را می‌توان روشی برای توصیف اثر نیروی الکتریکی در فضای اطراف بار دانست. اگر یک بار آزمایشی کوچک با اندازه‌ی \(q\) را در نقطه‌ای از میدان قرار دهیم و نیروی الکتریکی وارد بر آن را با \(F\) نشان دهیم، میدان الکتریکی در آن نقطه به‌صورت زیر تعریف می‌شود:

\[
\vec{E}=\frac{\vec{F}}{q}
\]

واحد میدان الکتریکی نیوتون بر کولن است که معادل ولت بر متر نیز هست. جهت میدان را با جهت نیروی وارد بر یک بار آزمایشی مثبت تعریف می‌کنیم. این میدان به ما کمک می‌کند بفهمیم بارهای الکتریکی چگونه در فضا بر یکدیگر اثر می‌گذارند، حتی اگر خود ذرات در تماس مستقیم نباشند.

در اینجا قصد نداریم وارد تمام جزئیات الکترومغناطیس شویم. آنچه برای ادامه لازم داریم این است که بار الکتریکی می‌تواند تحت تأثیر میدان الکتریکی نیرو دریافت کند و در شرایط مناسب حرکت کند. همین حرکت بارها، در قالبی مشخص، به چیزی منجر می‌شود که آن را جریان الکتریکی می‌نامیم.

\subsection{ولتاژ: اختلاف انرژی به ازای هر واحد بار}
\label{sub:voltage}

حالا به یکی از واژه‌هایی می‌رسیم که احتمالاً بیش از هر مفهوم دیگری در کار با برق شنیده‌اید: ولتاژ. ممکن است روی شارژر نوشته باشد ۵ ولت، روی باتری ۱۲ ولت، یا درباره‌ی یک مدار بگویند ولتاژ ورودی آن ۳٫۳ ولت است. اما این عدد دقیقاً چه چیزی را بیان می‌کند؟

برای فهم ولتاژ، بهتر است ابتدا انرژی را به یاد بیاوریم. وقتی یک جسم را از ارتفاع پایین به ارتفاع بالا می‌بریم، انرژی پتانسیل گرانشی آن تغییر می‌کند. در میدان الکتریکی نیز بارها می‌توانند انرژی پتانسیل الکتریکی داشته باشند. اگر بار الکتریکی را میان دو نقطه جابه‌جا کنیم، انرژی پتانسیل الکتریکی آن می‌تواند تغییر کند.

اختلاف پتانسیل الکتریکی میان دو نقطه را می‌توان تغییر انرژی پتانسیل الکتریکی به ازای واحد بار تعریف کرد:

\[
\Delta V=\frac{\Delta U}{q}
\]

در این رابطه، \(\Delta V\) اختلاف ولتاژ، \(\Delta U\) تغییر انرژی پتانسیل الکتریکی، و \(q\) مقدار بار است. واحد ولتاژ، ولت، برابر با یک ژول بر کولن است:

\[
1\;\mathrm{V}=1\;\mathrm{J/C}
\]

پس وقتی می‌گوییم میان دو نقطه اختلاف ولتاژ ۵ ولت وجود دارد، یعنی به ازای هر کولن بار، اختلاف انرژی پتانسیل الکتریکی متناظر با آن ۵ ژول است. این تعریف یک نکته‌ی مهم را روشن می‌کند: \textbf{ولتاژ خودِ انرژی نیست؛ انرژی به ازای واحد بار است.} \autocite{openstax-physics}

برای اینکه موضوع ملموس‌تر شود، می‌توانیم از یک تشبیه استفاده کنیم. در سامانه‌ی آبی، اختلاف ارتفاع دو مخزن می‌تواند زمینه‌ی حرکت آب را فراهم کند. در مدار الکتریکی نیز اختلاف پتانسیل میان دو نقطه می‌تواند زمینه‌ی حرکت بارها را فراهم کند. البته آب و بار الکتریکی یک چیز نیستند و این تشبیه فقط برای فهم ایده‌ی «اختلاف پتانسیل» مفید است؛ نباید آن را به‌عنوان مدل دقیق مدار در نظر بگیریم.

همچنین ولتاژ همیشه نسبت به یک مرجع سنجیده می‌شود. وقتی می‌گوییم نقطه‌ای ۵ ولت دارد، معمولاً منظورمان این است که نسبت به نقطه‌ی مرجع یا زمین مدار، اختلاف پتانسیل ۵ ولت دارد. اگر مرجع را عوض کنیم، مقدار بیان‌شده برای همان نقطه ممکن است عوض شود؛ آنچه برای عملکرد مدار اهمیت دارد، اختلاف پتانسیل میان نقاط موردنظر است.

\subsection{جریان الکتریکی: حرکت بار در زمان}
\label{sub:current}

ولتاژ به ما می‌گوید اختلاف پتانسیل الکتریکی میان دو نقطه چقدر است. اما وقتی بارهای الکتریکی در یک مسیر حرکت می‌کنند، با مفهوم دیگری روبه‌رو می‌شویم: جریان الکتریکی.

جریان الکتریکی نرخ عبور بار از یک مقطع مشخص در واحد زمان است. اگر در بازه‌ی زمانی \(\Delta t\)، مقدار بار \(\Delta Q\) از یک مقطع عبور کند، جریان متوسط برابر است با:

\[
I=\frac{\Delta Q}{\Delta t}
\]

و در بیان لحظه‌ای:

\[
I=\frac{dQ}{dt}
\]

واحد جریان، آمپر است و یک آمپر یعنی یک کولن بار در هر ثانیه از مقطع موردنظر عبور کند. \autocite{openstax-physics}

حالا می‌توانیم تفاوت ولتاژ و جریان را دقیق‌تر ببینیم. ولتاژ اختلاف پتانسیل میان دو نقطه است؛ جریان نرخ عبور بار است. وجود ولتاژ می‌تواند در شرایط مناسب باعث ایجاد جریان شود، اما این دو کمیت یکسان نیستند. ممکن است میان دو نقطه اختلاف ولتاژ وجود داشته باشد، اما به‌دلیل باز بودن مسیر، جریان پیوسته‌ای برقرار نشود. یا ممکن است در یک مدار، جریان برقرار باشد و هم‌زمان اختلاف ولتاژهایی میان اجزای مختلف آن وجود داشته باشد.

یک نکته‌ی تاریخی و قراردادی هم وجود دارد. جهت جریان متعارف را به‌عنوان جهت حرکت بار مثبت تعریف می‌کنیم. در بسیاری از رساناهای فلزی، حامل‌های بار منفی یعنی الکترون‌ها در جهت مخالف جریان متعارف حرکت می‌کنند. پس وقتی در نمودارهای مدار جهت جریان را با یک پیکان نشان می‌دهیم، لزوماً منظور جهت حرکت الکترون‌ها نیست؛ بلکه از قرارداد استاندارد جهت جریان استفاده می‌کنیم.

همچنین سرعت حرکت الکترون‌ها را نباید با سرعت انتشار تغییر سیگنال در مدار اشتباه بگیریم. در یک رسانای فلزی، حرکت رانشی متوسط الکترون‌ها ممکن است نسبتاً آهسته باشد، اما تغییر میدان الکتریکی که وضعیت مدار را تغییر می‌دهد، در امتداد مدار با سرعت بسیار بیشتری منتشر می‌شود. همین تفاوت است که اجازه می‌دهد تغییر وضعیت یک کلید یا یک ورودی، در مقیاس زمانی بسیار کوتاه در بخش‌های دیگر مدار اثر بگذارد.

\subsection{مقاومت، مدار و قانون اهم}
\label{sub:resistance-circuit-ohms-law}

برای اینکه جریان الکتریکی برقرار شود، معمولاً به یک مسیر رسانا و یک منبع اختلاف پتانسیل نیاز داریم. مسیر رسانا، اجزای مختلف و منبع می‌توانند یک مدار الکتریکی بسازند. مدار در ساده‌ترین حالت، مسیری است که امکان حرکت بار را فراهم می‌کند. اگر مسیر قطع باشد، جریان پیوسته‌ای برقرار نمی‌شود؛ اگر مسیر بسته باشد، بارها می‌توانند در مدار حرکت کنند.

اما همه‌ی مواد به یک شکل اجازه‌ی حرکت بار نمی‌دهند. مس، آلومینیوم و بسیاری از فلزها رساناهای خوبی هستند؛ درحالی‌که موادی مانند شیشه یا پلاستیک در شرایط معمول عایق‌اند. حتی در یک ماده‌ی رسانا نیز حرکت بار با محدودیت روبه‌رو است. این محدودیت را در سطحی ساده با کمیتی به نام مقاومت الکتریکی توصیف می‌کنیم.

مقاومت را با نماد \(R\) نشان می‌دهیم؛ واحد آن «اهم» نام دارد و نماد اختصاری آن حرف یونانی امگا است. برای یک عنصر اهمی در شرایطی که رابطه‌ی ولتاژ و جریان آن خطی است، قانون اهم به شکل زیر نوشته می‌شود:

\[
V=IR
\]

این رابطه می‌گوید اختلاف ولتاژ دو سر عنصر برابر حاصل‌ضرب جریان عبوری از آن در مقاومت آن است. برای مثال، اگر مقاومت یک عنصر ۱۰۰ اهم باشد و جریان ۰٫۰۱ آمپر از آن عبور کند، اختلاف ولتاژ دو سر آن ۱ ولت خواهد بود.

اما باید مراقب باشیم قانون اهم را به همه‌ی اجزای الکتریکی تعمیم ندهیم. بسیاری از اجزا در تمام محدوده‌های کاری خود رفتار خطی ندارند؛ دیودها، ترانزیستورها، لامپ‌های رشته‌ای و برخی مواد می‌توانند رابطه‌ی پیچیده‌تری میان ولتاژ و جریان داشته باشند. بنابراین قانون اهم برای عنصر اهمی و شرایط مناسب آن معتبر است، نه برای هر پدیده‌ی الکتریکی بدون استثنا. \autocite{openstax-physics}

در مدار، ولتاژ، جریان و مقاومت به یکدیگر وابسته‌اند؛ اما هرکدام چیزی جداگانه را توصیف می‌کنند. منبع می‌تواند اختلاف پتانسیل ایجاد کند، مسیر مدار امکان حرکت بار را فراهم می‌کند، و ویژگی اجزا تعیین می‌کند جریان چگونه برقرار شود. همین ارتباط‌ها پایه‌ی فهم مدارهای الکترونیکی‌اند، حتی اگر هنوز وارد طراحی مدارهای پیچیده نشده باشیم.

\subsection{اتصال سری و موازی: جریان از کدام مسیر می‌گذرد؟}
\label{sub:series-and-parallel-circuits}

تا اینجا مدار را به‌عنوان مسیری بسته برای حرکت بار شناختیم. اما در یک مدار واقعی، اجزا همیشه پشت سر هم و در یک مسیر واحد قرار نمی‌گیرند. گاهی جریان فقط یک مسیر برای عبور دارد؛ گاهی هم در یک نقطه میان چند شاخه تقسیم می‌شود. این تفاوت ساده، رفتار ولتاژ و جریان را تغییر می‌دهد و برای فهم مدارهای الکترونیکی اهمیت زیادی دارد.

وقتی دو یا چند جزء را به‌صورت سری متصل می‌کنیم، جریان عبوری از همه‌ی آن‌ها یکسان است؛ زیرا بار در یک مسیر پیوسته حرکت می‌کند و راه دیگری برای تقسیم شدن ندارد. در یک اتصال سری از مقاومت‌های ایده‌آل، مقاومت معادل برابر مجموع مقاومت‌هاست:

\[
R_{\mathrm{eq}}=R_1+R_2+\cdots+R_n
\]

در چنین مداری، اختلاف ولتاژ منبع میان اجزای مختلف تقسیم می‌شود. برای نمونه، اگر دو مقاومت ۱۰۰ و ۲۰۰ اهمی را به‌صورت سری به یک منبع ۹ ولتی ایده‌آل وصل کنیم، مقاومت کل ۳۰۰ اهم خواهد بود. با استفاده از قانون اهم، جریان مدار ۰٫۰۳ آمپر به دست می‌آید؛ بنابراین افت ولتاژ دو سر مقاومت اول ۳ ولت و دو سر مقاومت دوم ۶ ولت خواهد بود. مجموع این دو افت ولتاژ همان ۹ ولت منبع است.

در اتصال موازی، وضعیت برعکس است: دو سر هر شاخه به همان دو گره مشترک وصل می‌شوند و بنابراین اختلاف ولتاژ دو سر شاخه‌ها یکسان است. جریان کل میان شاخه‌ها تقسیم می‌شود و مقدار هر شاخه به مقاومت همان مسیر بستگی دارد. برای مقاومت‌های ایده‌آل، مقاومت معادل اتصال موازی از رابطه‌ی زیر به دست می‌آید:

\[
\frac{1}{R_{\mathrm{eq}}}=\frac{1}{R_1}+\frac{1}{R_2}+\cdots+\frac{1}{R_n}
\]

اگر همان مقاومت‌های ۱۰۰ و ۲۰۰ اهمی را این بار به‌صورت موازی به منبع ۹ ولتی وصل کنیم، هر دو مقاومت ۹ ولت دریافت می‌کنند؛ اما جریان شاخه‌ی ۱۰۰ اهمی ۰٫۰۹ آمپر و جریان شاخه‌ی ۲۰۰ اهمی ۰٫۰۴۵ آمپر خواهد بود. جریان منبع، در حالت ایده‌آل، مجموع جریان دو شاخه است. پس در اتصال سری، جریان یکسان می‌ماند و ولتاژ میان اجزا تقسیم می‌شود؛ در اتصال موازی، ولتاژ یکسان است و جریان میان شاخه‌ها تقسیم می‌شود.

این دو حالت فقط دو روش متفاوت برای کنار هم قرار دادن قطعات نیستند. آن‌ها نشان می‌دهند که «مدار» یک مجموعه‌ی بی‌ارتباط از اجزا نیست؛ آرایش اتصال‌ها تعیین می‌کند انرژی و بار چگونه در سامانه جابه‌جا شوند. در تحلیل دقیق‌تر مدارها از قوانین کیرشهف استفاده می‌کنیم: مجموع جریان‌های واردشونده به یک گره با مجموع جریان‌های خارج‌شونده برابر است، و مجموع تغییرات ولتاژ در یک حلقه‌ی بسته صفر می‌شود. این دو قانون، در سطح مدار، بازتابی از پایستگی بار و انرژی‌اند. \autocite{openstax-physics}

\subsection{توان الکتریکی: انرژی در واحد زمان}
\label{sub:electric-power}

تا اینجا از انرژی پتانسیل الکتریکی، ولتاژ و جریان حرف زدیم. اما اگر بخواهیم بدانیم یک جزء الکتریکی در هر ثانیه چه مقدار انرژی دریافت یا منتقل می‌کند، باید به مفهوم توان برسیم.

توان نرخ انتقال یا تبدیل انرژی در واحد زمان است. در یک مدار الکتریکی، برای یک عنصر دوپایه‌ای، توان لحظه‌ای به‌صورت زیر تعریف می‌شود:

\[
p(t)=v(t)i(t)
\]

در حالت جریان و ولتاژ مستقیم و پایدار، یا برای مقدارهای ثابت، همین رابطه به شکل ساده‌ی \(P=VI\) نوشته می‌شود.

اگر ولتاژ بر حسب ولت و جریان بر حسب آمپر باشد، توان بر حسب وات به دست می‌آید. برای مثال، اگر در یک لحظه اختلاف ولتاژ دو سر یک جزء ۵ ولت و جریان عبوری از آن ۰٫۲ آمپر باشد، توان متناظر برابر ۱ وات خواهد بود.

در یک مقاومت، بخشی از انرژی الکتریکی به گرما تبدیل می‌شود. اما در اجزای دیگر، انرژی می‌تواند به شکل نور، حرکت، میدان مغناطیسی یا تغییرات ذخیره‌شده در میدان الکتریکی و مغناطیسی ظاهر شود. به همین دلیل، توان فقط به معنای «گرمای تولیدشده» نیست؛ بلکه نرخ انتقال انرژی است. در یک مصرف‌کننده‌ی معمولی، انرژی الکتریکی ممکن است به شکل‌های دیگری تبدیل شود.

این موضوع برای بحث اطلاعات هم مهم است. وقتی یک سیگنال الکتریکی را تولید، تقویت، انتقال یا اندازه‌گیری می‌کنیم، در کنار نمایش اطلاعات، با محدودیت‌های انرژی و توان نیز روبه‌رو هستیم. یک سامانه‌ی واقعی نمی‌تواند بدون توجه به انرژی موردنیاز، گرما و محدودیت‌های فیزیکی، سیگنال را تا هر اندازه‌ای سریع یا قوی تغییر دهد.

\subsection{جریان مستقیم و متناوب: دو رفتار زمانی متفاوت}
\label{sub:dc-and-ac}

در گفت‌وگوهای روزمره احتمالاً با اصطلاح‌های جریان مستقیم و جریان متناوب روبه‌رو شده‌اید. باتری‌ها معمولاً منبع جریان مستقیم هستند که در منابع فنی با نماد \lr{DC} نشان داده می‌شود؛ درحالی‌که برق شبکه‌ی خانگی در بسیاری از کشورها به‌صورت جریان متناوب توزیع می‌شود که آن را با نماد \lr{AC} نشان می‌دهند. اما تفاوت این دو در چیست؟

در جریان مستقیم، جهت جریان در طول زمان ثابت می‌ماند؛ هرچند اندازه‌ی آن لزوماً ثابت نیست. پس \lr{DC} را نباید همیشه جریانی با مقدار کاملاً ثابت تصور کنیم. ممکن است شدت جریان بالا و پایین برود، اما جهت مرجع آن تغییر نکند.

در جریان متناوب، جریان به‌طور دوره‌ای جهت خود را عوض می‌کند. شکل موج آن در بسیاری از کاربردها تقریباً سینوسی است، هرچند شکل موج \lr{AC} الزاماً همیشه سینوسی نیست. برای یک موج سینوسی نمونه می‌توان نوشت:

\[
v(t)=A\sin(2\pi ft+\varphi)
\]

در این رابطه، \(A\) دامنه، \(f\) فرکانس، و \(\varphi\) فاز اولیه است. دامنه نشان می‌دهد تغییرات سیگنال تا چه اندازه از مقدار مرجع فاصله می‌گیرد؛ فرکانس تعداد چرخه‌ها در هر ثانیه را مشخص می‌کند؛ و فاز بیان می‌کند موج نسبت به یک مرجع زمانی در چه نقطه‌ای از چرخه قرار دارد.

این سه ویژگی در تحلیل سیگنال اهمیت دارند. دو سیگنال ممکن است فرکانس یکسان اما دامنه‌ی متفاوت داشته باشند؛ یا دامنه و فرکانس برابر اما فاز متفاوت داشته باشند. بسته به شیوه‌ی کدگذاری، هر کدام از این ویژگی‌ها می‌تواند برای انتقال اطلاعات به کار رود.

اما باید یک تمایز مهم را حفظ کنیم: جریان مستقیم و متناوب، به‌خودی‌خود معادل صفر و یک نیستند. \lr{AC} و \lr{DC} نوعی رفتار زمانی جریان یا ولتاژ را توصیف می‌کنند؛ صفر و یک نمادهای منطقی‌اند. یک سامانه‌ی دیجیتال می‌تواند با منبع \lr{DC} کار کند و درون آن سیگنال‌هایی داشته باشد که در زمان میان سطح‌های منطقی تغییر می‌کنند. بنابراین نباید بگوییم «\lr{DC} یعنی صفر» یا «\lr{AC} یعنی یک»؛ این‌ها به دو سطح متفاوت از توصیف تعلق دارند.

\subsection{خودِ سیگنال الکتریکی چه ویژگی‌هایی دارد؟}
\label{sub:electrical-signal-properties}

حالا که ولتاژ و جریان را می‌شناسیم، می‌توانیم سیگنال الکتریکی را دقیق‌تر توصیف کنیم. یک سیگنال الکتریکی ممکن است در طول زمان تغییر کند و همین تغییرات حامل پیام باشند. برای مطالعه‌ی آن، معمولاً به ویژگی‌هایی مانند دامنه، فرکانس، فاز، شکل موج و زمان بالا و پایین رفتن توجه می‌کنیم.

دامنه می‌تواند نشان دهد یک سیگنال تا چه حد از مقدار مرجع فاصله می‌گیرد. فرکانس مشخص می‌کند تغییرات تناوبی با چه سرعتی تکرار می‌شوند. فاز رابطه‌ی زمانی موج را با موج مرجع بیان می‌کند. شکل موج نیز می‌تواند سینوسی، مربعی، مثلثی یا نامنظم باشد. در یک مدار دیجیتال، سیگنال ممکن است مدت کوتاهی در یک سطح منطقی باقی بماند و سپس به سطح دیگر برود؛ در یک مدار آنالوگ، دامنه‌ی آن می‌تواند در یک بازه به‌صورت پیوسته تغییر کند.

برای مثال، اگر بخواهیم یک صدای سینوسی را با ولتاژ نمایش دهیم، ممکن است دامنه‌ی ولتاژ با فشار صوت متناظر شود. اما اگر بخواهیم یک بیت را در یک خط دیجیتال نشان دهیم، ممکن است ولتاژ در یک بازه‌ی مشخص به‌عنوان صفر و در بازه‌ای دیگر به‌عنوان یک تفسیر شود. در حالت اول، مقدار دامنه خودش بخشی از محتوای بازنمایی‌شده است؛ در حالت دوم، معمولاً اهمیت اصلی با این است که ولتاژ در کدام محدوده‌ی معتبر قرار گرفته است.

در یک سامانه‌ی دیجیتال، گیرنده سیگنال را اندازه‌گیری می‌کند و با توجه به آستانه‌های تعریف‌شده، تصمیم می‌گیرد کدام نماد دریافت شده است. این تصمیم‌گیری همیشه با محدودیت دقت و نویز همراه است. اگر سیگنال به‌اندازه‌ی کافی از آستانه فاصله داشته باشد، تغییرهای کوچک ممکن است اثری بر نتیجه نگذارند؛ اما اگر در نزدیکی مرز تصمیم‌گیری قرار بگیرد، احتمال اشتباه بیشتر می‌شود.

به همین دلیل است که در طراحی سامانه‌های دیجیتال، فقط مقدار اسمی صفر و یک کافی نیست. باید محدوده‌ی ولتاژ قابل قبول، زمان پایدار ماندن سیگنال، سرعت تغییر آن و میزان نویزی را که سامانه می‌تواند تحمل کند نیز مشخص کنیم. اطلاعات زمانی که واقعاً قابل اعتماد منتقل می‌شوند، نه فقط به کدگذاری منطقی، بلکه به رفتار فیزیکی سیگنال وابسته‌اند.

\subsection{خازن و سلف: دو راه برای ذخیره‌ی انرژی در میدان‌ها}
\label{sub:capacitors-and-inductors}

پیش از پایان این فصل، بد نیست با دو جزء مهم دیگر در مدارها آشنا شویم: خازن و سلف. در اینجا قرار نیست طراحی یا تحلیل کامل مدارهای دارای این اجزا را انجام دهیم؛ هدف فقط این است که ببینیم چرا رفتار سیگنال در مدار، به اجزای فیزیکی آن وابسته است.

خازن قطعه‌ای است که می‌تواند انرژی را در میدان الکتریکی ذخیره کند. در ساده‌ترین مدل، دو رسانا را در نظر بگیرید که با ماده‌ای عایق از یکدیگر جدا شده‌اند. وقتی میان آن‌ها اختلاف ولتاژ ایجاد می‌شود، بارهای الکتریکی روی دو رسانا با علامت‌های مخالف جمع می‌شوند و میدان الکتریکی در فضای میان آن‌ها شکل می‌گیرد. ظرفیت خازن با نماد \(C\) نشان داده می‌شود و برای یک خازن ایده‌آل رابطه‌ی زیر را داریم:

\[
Q=CV
\]

یعنی مقدار بار ذخیره‌شده با ظرفیت و ولتاژ متناسب است. انرژی ذخیره‌شده در خازن ایده‌آل نیز برابر است با:

\[
U_C=\frac{1}{2}CV^2
\]

خازن در برابر تغییر ناگهانی ولتاژ رفتاری خاص نشان می‌دهد، چون برای تغییر ولتاژ دو سر آن باید بار ذخیره‌شده تغییر کند. به همین دلیل، خازن‌ها می‌توانند در صاف‌کردن نوسان‌های ولتاژ، فیلتر کردن سیگنال‌ها، ایجاد تأخیرهای زمانی و بسیاری از کاربردهای دیگر نقش داشته باشند.

سلف یا القاگر، در مقابل، انرژی را در میدان مغناطیسی ذخیره می‌کند. هنگامی که جریان از سیم‌پیچی عبور می‌کند، میدان مغناطیسی ایجاد می‌شود. در مدل ایده‌آل، ولتاژ دو سر سلف با نرخ تغییر جریان متناسب است:

\[
v_L=L\frac{di}{dt}
\]

در این رابطه، \(L\) اندوکتانس یا ضریب القای سلف است. بنابراین سلف در برابر تغییر سریع جریان مقاومت می‌کند؛ نه به این معنا که مانند یک مقاومت معمولی انرژی را به گرما تبدیل می‌کند، بلکه به این دلیل که تغییر جریان با تغییر میدان مغناطیسی و ولتاژ القایی همراه است. انرژی ذخیره‌شده در یک سلف ایده‌آل برابر است با:

\[
U_L=\frac{1}{2}LI^2
\]

این دو جزء یادآوری می‌کنند که مدار فقط مجموعه‌ای از سیم‌ها و عددهای ولتاژ نیست. در مدار واقعی، انرژی می‌تواند در میدان‌های الکتریکی و مغناطیسی ذخیره شود، میان اجزا جابه‌جا شود و در طول زمان تغییر کند. همین ویژگی‌ها شکل و زمان‌بندی سیگنال را تعیین می‌کنند. از این رو، وقتی درباره‌ی سیگنال الکتریکی حرف می‌زنیم، باید هم به پیام منطقی و هم به رفتار فیزیکی اجزایی توجه کنیم که آن سیگنال را تولید و منتقل می‌کنند.

\subsection{یک مثال کامل: از دمای اتاق تا بیت}
\label{sub:temperature-to-bit}

تا اینجا مفاهیم را جداگانه بررسی کردیم. حالا بیایید همه‌ی آن‌ها را در یک مثال کنار هم بگذاریم. فرض کنید یک دستگاه اندازه‌گیری دما داریم که باید دمای اتاق را اندازه بگیرد، مقدار آن را به یک رایانه بدهد و روی صفحه نمایش دهد.

در جهان واقعی، دمای اتاق یک کمیت فیزیکی است که در طول زمان تغییر می‌کند. حسگر دما به تغییر دما واکنش نشان می‌دهد و در خروجی خود یک کمیت الکتریکی متناسب با آن تولید می‌کند. در یک مدل ساده، ممکن است خروجی حسگر ولتاژی باشد که با دما تغییر می‌کند. در این نقطه، با یک سیگنال آنالوگ روبه‌رو هستیم.

اما رایانه نمی‌تواند صرفاً به یک مفهوم انتزاعی به نام «دمای اتاق» دسترسی داشته باشد. باید مقدار ولتاژ را اندازه بگیرد. مبدل آنالوگ به دیجیتال، که به‌اختصار \lr{ADC} نامیده می‌شود، این کار را انجام می‌دهد: در زمان‌های مشخص از سیگنال نمونه می‌گیرد و سپس مقدار هر نمونه را به یکی از سطح‌های دیجیتال نگاشت می‌کند. هر سطح با یک کد دودویی نمایش داده می‌شود.

فرض کنید مبدل ما ۱۲ بیت دقت دارد. در این صورت برای هر نمونه می‌تواند یکی از \(2^{12}=4096\) کد متفاوت را تولید کند. این کد نماینده‌ی یک محدوده‌ی مشخص از مقدار ورودی است. هرچه تعداد بیت‌ها بیشتر باشد، تعداد سطح‌های قابل نمایش افزایش پیدا می‌کند؛ اما دقت نهایی همچنان به نویز حسگر، کیفیت مدار، مرجع ولتاژ، طراحی مبدل و کالیبراسیون دستگاه وابسته است.

پس از تبدیل، رایانه رشته‌ای از بیت‌ها را دریافت می‌کند. این بیت‌ها در سطح منطقی عددهایی هستند که می‌توان آن‌ها را ذخیره، پردازش یا منتقل کرد. اما در داخل دستگاه، این صفر و یک‌ها با وضعیت‌های الکتریکی مشخص نمایش داده می‌شوند. در نهایت، نرم‌افزار مقدار دیجیتال را طبق مقیاس و واحد حسگر تفسیر می‌کند و مثلاً عدد ۲۳٫۵ درجه‌ی سانتی‌گراد را روی نمایشگر نشان می‌دهد.

اگر بخواهیم کل مسیر را خلاصه کنیم، چنین زنجیره‌ای داریم:

\begin{figure}[htbp]
\centering
\begin{tikzpicture}[node distance=0.36cm, every node/.style={draw,rounded corners,align=center,minimum width=5.2cm,minimum height=0.72cm}, >=Latex]
  \node (temp) {دمای واقعی اتاق};
  \node[below=of temp] (sensor) {حسگر و تبدیل دما به ولتاژ};
  \node[below=of sensor] (analog) {سیگنال آنالوگ};
  \node[below=of analog] (adc) {نمونه‌برداری و کوانتیزه‌سازی};
  \node[below=of adc] (bits) {کد دودویی};
  \node[below=of bits] (screen) {تفسیر نرم‌افزاری و نمایش عدد};
  \draw[->,thick] (temp) -- (sensor);
  \draw[->,thick] (sensor) -- (analog);
  \draw[->,thick] (analog) -- (adc);
  \draw[->,thick] (adc) -- (bits);
  \draw[->,thick] (bits) -- (screen);
\end{tikzpicture}
\caption{نمونه‌ای از زنجیره‌ی تبدیل یک کمیت فیزیکی به داده‌ی دیجیتال و بازنمایی دوباره‌ی آن برای انسان.}
\label{fig:temperature-to-bit}
\end{figure}

در این مثال، دما، ولتاژ، نمونه‌های عددی، بیت‌ها و عدد نمایش‌داده‌شده روی صفحه، پنج چیز کاملاً یکسان نیستند. هر کدام در سطحی متفاوت از فرایند قرار دارند و رابطه‌ای مشخص میان آن‌ها برقرار است. ما یک کمیت فیزیکی را اندازه می‌گیریم، به سیگنال تبدیل می‌کنیم، آن را به نمادهای گسسته نگاشت می‌کنیم و سپس دوباره از طریق نرم‌افزار و نمایشگر به شکلی قابل فهم برای انسان بازمی‌گردانیم.

این مثال به ما نشان می‌دهد که داده در جهان محاسباتی، صرفاً یک عدد یا رشته‌ی صفر و یک نیست. در پشت آن، زنجیره‌ای از اندازه‌گیری، بازنمایی، انتقال، تشخیص و تفسیر وجود دارد. هرجا که یکی از این پیوندها خراب شود، ممکن است داده‌ی نهایی با واقعیت اولیه تفاوت پیدا کند. اگر حسگر خطا داشته باشد، ولتاژ نادرست تولید می‌شود؛ اگر نویز زیاد باشد، مقدار اشتباه اندازه‌گیری می‌شود؛ اگر نمونه‌برداری مناسب نباشد، تغییرات سریع از دست می‌روند؛ و اگر کدگذاری یا تفسیر نادرست باشد، عدد نهایی معنای اشتباهی پیدا می‌کند.

\section{از بازنمایی تا واقعیت فیزیکی}
\label{sec:from-representation-to-physical-reality}

در این فصل از پرسشی شروع کردیم که در پایان فصل قبل باقی مانده بود: یک بازنمایی انتزاعی چگونه در جهان واقعی وجود پیدا می‌کند؟ دیدیم که برای پاسخ دادن به آن، ابتدا باید اطلاعات را از معنا جدا کنیم و بفهمیم نظریه‌ی شانون دقیقاً چه چیزی را اندازه می‌گیرد. سپس با احتمال، عدم‌قطعیت، آنتروپی، نویز و ظرفیت کانال آشنا شدیم و دیدیم که پیچیدگی کولموگروف پرسشی متفاوت از آنتروپی شانون دارد.

بعد سراغ خودِ بازنمایی رفتیم و دیدیم یک پیام می‌تواند از طریق حامل‌های فیزیکی متفاوتی نمایش داده شود. با سه‌گانه‌ی پیرس آشنا شدیم تا بفهمیم یک نشانه چگونه از راه شباهت، پیوند واقعی یا قرارداد به موضوع خود مربوط می‌شود. همچنین میان پیام، کد، سیگنال و کانال تفاوت گذاشتیم؛ چهار مفهومی که در گفت‌وگوی روزمره گاهی به‌جای هم به کار می‌روند، اما در تحلیل یک سامانه‌ی اطلاعاتی نقش‌های متفاوتی دارند.

در ادامه دیدیم که اطلاعات و حامل آن یکی نیستند، هرچند پردازش و ذخیره‌سازی اطلاعات در جهان واقعی به وضعیت‌های فیزیکی وابسته است. بحث شیطان ماکسول و اصل لاندائر نشان داد که در برخی عملیات برگشت‌ناپذیر، مانند پاک کردن یک بیت نامعلوم، محدودیت‌های ترمودینامیکی اهمیت پیدا می‌کنند. البته این اصل را نباید به تمام عملیات محاسباتی تعمیم دهیم؛ آنچه درباره‌ی هزینه‌ی فیزیکی گفتیم، به شرایط و نوع عملیات مشخصی مربوط است.

سپس از سیگنال پیوسته به نمایش گسسته حرکت کردیم. نمونه‌برداری، کوانتیزه‌سازی و کدگذاری نشان دادند که یک پدیده‌ی پیوسته چگونه می‌تواند به مجموعه‌ای از نمادهای محدود تبدیل شود. با وجود این، سیگنال دیجیتال نیز همچنان یک پدیده‌ی فیزیکی است؛ فقط گیرنده به‌جای تفسیر بی‌نهایت مقدار ممکن، سطح‌ها و محدوده‌های مشخصی را به نمادهای منطقی نسبت می‌دهد.

و در پایان، مسیر را با الکتریسیته ادامه دادیم. بار الکتریکی، میدان، ولتاژ، جریان، مقاومت و توان را شناختیم و دیدیم چرا اختلاف پتانسیل و جریان می‌توانند در مدارهای الکترونیکی برای نمایش و انتقال اطلاعات به کار روند. خازن و سلف نیز نشان دادند که حتی یک سیگنال ساده، در عمل با انرژی ذخیره‌شده در میدان‌های الکتریکی و مغناطیسی و با رفتار اجزای واقعی مدار پیوند دارد.

پس ایده‌ی اصلی این فصل را می‌توان در یک جمله خلاصه کرد: \textbf{اطلاعات هرچند به‌صورت انتزاعی توصیف می‌شود، برای ثبت، انتقال یا پردازش باید در قالب حالت‌های فیزیکیِ قابل تشخیص بازنمایی شود.}
